\documentclass[11pt,a4paper]{ctexart}

\usepackage{fontspec}
\IfFontExistsTF{Times New Roman}{
  \setmainfont{Times New Roman}
}{
  \setmainfont{Liberation Serif}
}
\IfFontExistsTF{SimSun}{
  \setCJKmainfont[AutoFakeBold=2.5]{SimSun}
}{
  \IfFontExistsTF{Songti SC}{
    \setCJKmainfont{Songti SC}
  }{
    \IfFontExistsTF{STSong}{
      \setCJKmainfont{STSong}
    }{
      \setCJKmainfont{Noto Serif CJK SC}
    }
  }
}

\usepackage[margin=0.85in]{geometry}
\usepackage{amsmath,amssymb,amsthm,mathtools}
\usepackage{booktabs,array,longtable}
\usepackage{xcolor}
\usepackage{tcolorbox}
\tcbuselibrary{skins,breakable}
\usepackage{enumitem}
\usepackage{microtype}
\usepackage{titlesec}
\usepackage{fancyhdr}
\usepackage[round,authoryear]{natbib}
\usepackage[colorlinks=true,linkcolor=teal!80!black,citecolor=blue!80!black,urlcolor=teal!80!black]{hyperref}

\pagestyle{fancy}
\fancyhf{}
\fancyhead[L]{\small\color{gray}理论推导手册：环境方差与报告异质性下再生数估计}
\fancyhead[R]{\small\color{gray}王驰 $\cdot$ 西安科技大学}
\fancyfoot[C]{\thepage}
\renewcommand{\headrulewidth}{0.4pt}

% 数学记号宏定义
\newcommand{\E}{\mathbb{E}}
\newcommand{\Var}{\mathrm{Var}}
\newcommand{\Cov}{\mathrm{Cov}}
\newcommand{\NB}{\mathrm{NB}}
\newcommand{\Bin}{\mathrm{Bin}}
\newcommand{\Beta}{\mathrm{Beta}}
\newcommand{\CV}{\mathrm{CV}}
\newcommand{\Rbar}{\bar{R}}
\newcommand{\mx}{m_\times}
\newcommand{\mxH}{m_{\times,H}}

\newtheorem{definition}{定义}[section]
\newtheorem{theorem}{定理}[section]
\newtheorem{proposition}{命题}[section]
\newtheorem{corollary}{推论}[section]
\newtheorem{remark}{注}[section]

% 自定义高质感彩色方框
\newtcolorbox{theorembox}[1]{
  colback=blue!3!white,
  colframe=blue!60!black,
  fonttitle=\bfseries\sffamily,
  title={#1},
  arc=2mm,
  boxrule=0.8pt,
  breakable,
  enhanced,
  drop shadow=black!10
}

\newtcolorbox{derivationbox}[1]{
  colback=gray!4!white,
  colframe=gray!60!black,
  fonttitle=\bfseries\sffamily,
  title={#1},
  arc=1.5mm,
  boxrule=0.6pt,
  breakable,
  enhanced
}

\newtcolorbox{insightbox}[1]{
  colback=teal!4!white,
  colframe=teal!70!black,
  fonttitle=\bfseries\sffamily,
  title={#1},
  arc=1.5mm,
  boxrule=0.6pt,
  breakable,
  enhanced
}

\title{\Huge\bfseries\color{blue!80!black} 环境方差与报告异质性下再生数估计的\\完整数学推导手册\\[12pt]
\Large\normalfont\color{gray!80!black} 从微观分支过程、信息论界限到监测设计预算优化的分步推导与证明}
\author{\large\textbf{王驰}\\[4pt]
\normalsize 西安科技大学 安全科学与工程学院\\[2pt]
\small\texttt{wangc@bostech.com.cn}}
\date{\small 2026年10月 $\cdot$ 完整推导版}

\begin{document}

\maketitle
\thispagestyle{empty}

\begin{abstract}
\noindent 本手册针对论文《环境方差与报告异质性下再生数估计的不可约误差与监测设计》（及英文手稿 \emph{Limits to improving reproduction number estimates by increasing case counts: Environmental variability, reporting heterogeneity, and surveillance design}）提供完备、自洽、详尽到每一步代数展开的数学推导与理论证明。手册从负二项分支过程与二项稀疏的概率母函数（PGF）保持定理出发，严格区分条件瞬时再生数 $R_t$（估计目标 1）与跨期平均再生数 $\bar{R}$（估计目标 2）；详细梳理了条件无偏估计类中 Cramér--Rao 下界的达到性、全方差公式展开下的环境方差底板 $v_R + A/m$、方差平衡规模 $m_\times$ 的物理几何性质、报告异质性引入的单期二阶矩简并边界与诊断性报告签名 $F(w, \psi)$、覆盖度与报告率的决策反转边界 $m > \rho^2\nu/k$、一阶自回归与空间汇总下的预算阈值 $\mathcal{B}^\dagger$、环境方差估计量 $\hat{v}_R$ 的渐近平台 $\Omega_\infty/T$ 与季节相对误差下界 $\sqrt{2/S}$、动态更新过程中的二阶 Delta 展开以及线性收缩经验贝叶斯估计的方差—偏差权衡。手册旨在为学术报告、理论答辩及同行评议提供清晰、规范的数学分析底本。
\end{abstract}

\vspace{10pt}
\tableofcontents
\newpage

\section{记号系统、概率分布与核心假设}
\label{sec:part01}

\subsection{数学符号全览}

为了使推导过程具备完全的自洽性与可溯源性，下表系统总结了本推导手册及论文全文所使用的全部核心数学符号。

\begin{center}
\small
\begin{longtable}{@{}llll@{}}
\toprule
\textbf{符号} & \textbf{数学含义} & \textbf{量纲 / 典型取值} & \textbf{出现小节} \\
\midrule
\endfirsthead
\toprule
\textbf{符号} & \textbf{数学含义} & \textbf{量纲 / 典型取值} & \textbf{出现小节} \\
\midrule
\endhead
\bottomrule
\endfoot
$t$ & 离散世代传播时期（周或代） & 时间（整数） & \S\ref{sec:part02} \\
$n$ & 当期上一代有效感染者人数（祖代规模） & 人数（整数） & \S\ref{sec:part02} \\
$X_{i,t}$ & 第 $t$ 期第 $i$ 名感染者产生的继发感染数 & 人数（非负整数） & \S\ref{sec:part02} \\
$R_t$ & 第 $t$ 期的条件瞬时再生数（$X_{i,t}$ 的条件期望） & 无量纲，$[0, \infty)$ & \S\ref{sec:part02} \\
$k$ & 负二项超传播离散参数（Dispersion parameter） & 无量纲，流感 $\sim 0.35$ & \S\ref{sec:part02} \\
$I_t$ & 第 $t$ 期产生的新发感染总数（$\sum_{i=1}^n X_{i,t}$） & 人数（非负整数） & \S\ref{sec:part02} \\
$\rho$ & 预期报告概率（Reporting completeness） & 概率，$[0, 1]$ & \S\ref{sec:part02} \\
$C_t$ & 第 $t$ 期监测捕获的报告确诊病例数 & 例数（非负整数） & \S\ref{sec:part02} \\
$m$ & 预期报告规模（Expected reporting scale，$m \equiv \rho n$） & 例数 / 期 & \S\ref{sec:part02} \\
$\widehat{R}_t$ & 自然估计量（Natural plug-in estimator，$C_t / m$） & 无量纲 & \S\ref{sec:part03} \\
$\Rbar$ & 跨期边际平均再生数（$\E(R_t)$） & 无量纲，$\sim 1.15$ 或 $1.67$ & \S\ref{sec:part04} \\
$v_R$ & 跨期环境方差（Environmental variance，$\Var(R_t)$） & 无量纲，$\sim 0.04$ & \S\ref{sec:part04} \\
$A$ & 边际方差规模项系数（Scale coefficient） & 无量纲 & \S\ref{sec:part04} \\
$\mx$ & 方差平衡规模（Variance balance scale，$A / v_R$） & 例数，基准 $\sim 53.1$ & \S\ref{sec:part05} \\
$p_t$ & 第 $t$ 期实际发生的随机报告概率 & 概率，$[0, 1]$ & \S\ref{sec:part06} \\
$\tau$ & 报告率的相对变异系数（$\sigma / \rho$） & 无量纲，基准 $\sim 0.15$ & \S\ref{sec:part06} \\
$u_t$ & 标准化报告率（$p_t / \rho$，$\E(u_t)=1$） & 无量纲 & \S\ref{sec:part06} \\
$v_H$ & 报告异质性方差（Reporting variance，$\tau^2(\Rbar^2+v_R)$） & 无量纲 & \S\ref{sec:part06} \\
$A_H$ & 计入报告波动后的修正规模项系数 & 无量纲 & \S\ref{sec:part06} \\
$\mxH$ & 报告异质性修正后的平衡规模 & 例数，基准 $\sim 30.4$ & \S\ref{sec:part06} \\
$U$ & 统计可识别复合参数（$(1+\tau^2)(\Rbar^2+v_R)$） & 无量纲 & \S\ref{sec:part06} \\
$\nu$ & 二项型报告抽样的先验样本量（Beta-Binomial 参数） & 样本数 & \S\ref{sec:part07} \\
$D_t$ & 加权回溯分母规模（Surveillance denominator） & 例数 / 周 & \S\ref{sec:part08} \\
$G_t$ & 经验周增长比统计量（$C_{t+1}/D_t$） & 无量纲 & \S\ref{sec:part08} \\
$w_s$ & 世代间隔或回溯加权分布权重 & 权重向量，$\sum w_s = 1$ & \S\ref{sec:part08} \\
$\psi$ & 报告冲击的周间一阶自相关系数 & 相关系数，$[-1, 1]$ & \S\ref{sec:part08} \\
$F(w, \psi)$ & 报告签名卷积核函数 & 无量纲 & \S\ref{sec:part08} \\
$q_1, q_2$ & 改变分母回溯权重引起的截距比值 & 无量纲比值 & \S\ref{sec:part08} \\
$T$ & 连续观测的总期数 & 周期数 & \S\ref{sec:part09} \\
$\phi$ & 环境波动 $R_t$ 的时间一阶自相关系数 & 相关系数，$[-1, 1]$ & \S\ref{sec:part09} \\
$T_{\mathrm{eff}}$ & 方差估计的有效观测期数 & 周期数 & \S\ref{sec:part09} \\
$\mathcal{B}$ & 监测总财政/资源预算 & 资金 / 成本单位 & \S\ref{sec:part09} \\
$c_0$ & 每多监测一期的固定运维成本 & 成本单位 / 期 & \S\ref{sec:part09} \\
$c_1$ & 每多追踪捕获一例报告病例的边际成本 & 成本单位 / 例 & \S\ref{sec:part09} \\
$\mathcal{B}^\dagger$ & 自相关环境下规模/期数决策反转的关键预算阈值 & 成本单位 & \S\ref{sec:part09} \\
$J$ & 空间协同监测的辖区（州/市）数量 & 辖区数 & \S\ref{sec:part09} \\
$\bar{c}$ & 空间辖区间环境扰动的平均两两相关系数 & 相关系数，$[-1, 1]$ & \S\ref{sec:part09} \\
$\hat{v}_R$ & 环境方差的无偏修正矩估计量 & 无量纲 & \S\ref{sec:part10} \\
$\Omega_\infty$ & 环境方差估计方差在规模无穷大时的极限平台 & 无量纲 & \S\ref{sec:part10} \\
$S$ & 观测到的独立流行季（Epidemic seasons）数目 & 季节数 & \S\ref{sec:part10} \\
$v_S$ & 流行季层面的共同随机效应方差 & 无量纲 & \S\ref{sec:part10} \\
$\beta(n, \delta)$ & 检出 $R_t=1+\delta$ 暴发增长的统计功效（Power） & 概率，$[0, 1]$ & \S\ref{sec:part11} \\
$\Lambda_t$ & 更新过程中的真实感染压力（Infection pressure） & 人数 & \S\ref{sec:part12} \\
$\CV(\Lambda)$ & 感染压力的变异系数 & 无量纲 & \S\ref{sec:part12} \\
$\lambda$ & 线性收缩估计量的收缩权重系数 & 无量纲，$[0, 1]$ & \S\ref{sec:part13} \\
\bottomrule
\end{longtable}
\end{center}

\subsection{负二项分布的标准参数化约定}

在统计物理、生态学与流行病学中，负二项分布存在多种不同的参数化方式。本手册严格采用概率母函数（PGF）在分支过程中最自然的“失败次数（Number of failures）”参数化约定：
\begin{definition}[负二项分布 $\NB(r, p)$]
设随机变量 $Y \sim \NB(r, p)$，其中 $r > 0$ 为正实数（分散参数），$p \in (0, 1]$ 为单次伯努利试验的成功概率。其概率质量函数（PMF）为：
\begin{equation}
  P(Y = y) = \frac{\Gamma(y + r)}{y! \, \Gamma(r)} \, p^r (1 - p)^y, \qquad y = 0, 1, 2, \dots
  \label{eq:nb-pmf}
\end{equation}
其概率母函数定义为 $G_Y(s) \equiv \E(s^Y)$：
\begin{align}
  G_Y(s) &= \sum_{y=0}^\infty s^y \frac{\Gamma(y+r)}{y!\,\Gamma(r)} p^r (1-p)^y 
  = p^r \sum_{y=0}^\infty \binom{-r}{y} (-1)^y [(1-p)s]^y \nonumber \\
  &= p^r \big(1 - (1-p)s\big)^{-r} = \left( \frac{p}{1 - (1-p)s} \right)^r.
  \label{eq:nb-pgf}
\end{align}
\end{definition}

利用概率母函数的求导性质，可立即导出其理论矩：
\begin{itemize}
  \item 一阶导数与数学期望：
  \begin{equation}
    G_Y'(s) = r \left( \frac{p}{1 - (1-p)s} \right)^{r-1} \frac{p(1-p)}{\big(1-(1-p)s\big)^2} = \frac{r(1-p)}{p} \left( \frac{p}{1 - (1-p)s} \right)^{r+1},
  \end{equation}
  令 $s = 1$，得到数学期望：
  \begin{equation}
    \E(Y) = G_Y'(1) = \frac{r(1-p)}{p}.
    \label{eq:nb-mean}
  \end{equation}
  \item 二阶导数与方差：
  \begin{equation}
    G_Y''(s) = \frac{r(r+1)(1-p)^2}{p^2} \left( \frac{p}{1 - (1-p)s} \right)^{r+2},
  \end{equation}
  令 $s = 1$，得到二阶阶乘矩 $\E[Y(Y-1)] = G_Y''(1) = \frac{r(r+1)(1-p)^2}{p^2}$。由方差公式：
  \begin{align}
    \Var(Y) &= \E[Y(Y-1)] + \E(Y) - [\E(Y)]^2 \nonumber \\
    &= \frac{r(r+1)(1-p)^2}{p^2} + \frac{r(1-p)}{p} - \frac{r^2(1-p)^2}{p^2} \nonumber \\
    &= \frac{r(1-p)^2}{p^2} + \frac{r(1-p)}{p} = \frac{r(1-p)}{p^2}.
    \label{eq:nb-var}
  \end{align}
\end{itemize}

\subsection{公理化核心假设体系（H1--H6）}

推导基于以下六条统计与流行病学假设。前四条界定了理想单期监测基准，后两条拓展至动态现实场景：

\begin{theorembox}{核心假设公理系统}
\begin{description}
  \item[H1（负二项分支与独立二项观测）：] 给定当期繁殖条件 $R_t$，各感染者的继发感染数相互独立，且服从参数为 $r=k, p=k/(k+R_t)$ 的负二项分布；各新发感染者独立地以概率 $\rho \in (0, 1]$ 被监测系统捕获上报。
  \item[H2（同世代环境共享）：] 处于同一传播世代 $t$ 内的全体感染个体共享同一个外生再生数 $R_t$；环境随机性作用于“世代”或“周”整体，而非个体层面。
  \item[H3（外生参数已知性）：] 超传播过离散参数 $k$ 与平均报告率 $\rho$ 属于病原体生物学与监测系统的固有外生参数，在推导基准中视为已知。
  \item[H4（祖代规模的外生独立性）：] 上一代感染规模 $n$ 由前瞻性调查设计（如主动接触者追踪）固定，统计独立于当期环境随机扰动，即满足：
  \begin{equation}
    \E(R_t \mid n) = \Rbar, \qquad \Var(R_t \mid n) = v_R.
  \end{equation}
  对于一步前瞻预测，进一步要求 $R_{t+1}$ 独立于截至第 $t$ 期的全部历史观测信息集 $\mathcal{F}_t$。
  \item[H5（随机报告率过程）：] 报告率 $p_t \in (0, 1]$ 构成平稳随机过程，与真实的传播动力学过程 $\{R_t\}$ 及祖代规模 $n$ 相互独立。其一阶矩与方差为 $\E(p_t) = \rho$，$\Var(p_t) = \sigma^2$。定义相对报告波动率为 $\tau \equiv \sigma / \rho$，标准化报告率为 $u_t \equiv p_t / \rho$（满足 $\E(u_t) = 1, \Var(u_t) = \tau^2$）。
  \item[H6（多期条件独立性）：] 在给定环境扰动轨道 $(R_1, \dots, R_T)$ 与报告率轨道 $(p_1, \dots, p_T)$ 的条件下，各期的报告病例生成过程相互独立。
\end{description}
\end{theorembox}

\section{微观负二项分支过程与二项稀疏定理}
\label{sec:part02}

\subsection{超传播个体的泊松--伽马混合机制}

在流行病学中，负二项后代分布被广泛用于刻画“超级传播（Superspreading）”现象 \citep{lloydsmith2005superspreading}。本小节从连续型个体传染力的伽马混合出发，显式推导式~\eqref{eq:offspring} 的负二项概率质量函数与理论矩。

\begin{derivationbox}{机制推导 2.1：泊松--伽马混合生成负二项分布}
设感染个体 $i$ 的内在传染力（传代强度）为随机变量 $\nu_i$。由于个体免疫状态、接触网络度数及微环境暴露的异质性，假设 $\nu_i$ 服从均值为 $R_t$、形状参数为 $k$ 的伽马（Gamma）分布：
\begin{equation}
  \nu_i \sim \text{Gamma}\!\left(k,\ \frac{k}{R_t}\right), \qquad f_{\nu}(\nu) = \frac{(k/R_t)^k}{\Gamma(k)} \, \nu^{k-1} e^{-\frac{k}{R_t}\nu}, \quad \nu > 0.
\end{equation}
给定该传染力 $\nu_i$，其产生的继发感染数 $X_i$ 服从均值为 $\nu_i$ 的泊松分布：
\begin{equation}
  X_i \mid \nu_i \sim \text{Poisson}(\nu_i), \qquad P(X_i = x \mid \nu_i) = \frac{e^{-\nu_i} \nu_i^x}{x!}, \quad x = 0, 1, 2, \dots
\end{equation}
对潜变量 $\nu_i$ 积分求边际分布：
\begin{align}
  P(X_i = x \mid R_t) &= \int_0^\infty P(X_i = x \mid \nu) f_\nu(\nu) \, \mathrm{d}\nu \nonumber \\
  &= \int_0^\infty \frac{e^{-\nu} \nu^x}{x!} \frac{(k/R_t)^k}{\Gamma(k)} \nu^{k-1} e^{-\frac{k}{R_t}\nu} \, \mathrm{d}\nu \nonumber \\
  &= \frac{(k/R_t)^k}{x! \, \Gamma(k)} \int_0^\infty \nu^{x+k-1} e^{-\left(1 + \frac{k}{R_t}\right)\nu} \, \mathrm{d}\nu.
\end{align}
引入变量代换 $u = \left(1 + \frac{k}{R_t}\right)\nu = \frac{k+R_t}{R_t}\nu$，利用伽马函数定义 $\int_0^\infty u^{x+k-1} e^{-u} \mathrm{d}u = \Gamma(x+k)$：
\begin{align}
  P(X_i = x \mid R_t) &= \frac{(k/R_t)^k}{x! \, \Gamma(k)} \left( \frac{R_t}{k+R_t} \right)^{x+k} \Gamma(x+k) \nonumber \\
  &= \frac{\Gamma(x+k)}{x! \, \Gamma(k)} \left( \frac{k}{k+R_t} \right)^k \left( \frac{R_t}{k+R_t} \right)^x.
  \label{eq:der-nb-mix}
\end{align}
对照负二项定义式~\eqref{eq:nb-pmf}，式~\eqref{eq:der-nb-mix} 严格等价于参数为 $r = k$ 且成功概率为
\begin{equation}
  p = \frac{k}{k+R_t}, \qquad 1 - p = \frac{R_t}{k+R_t}
\end{equation}
的负二项分布 $\NB\big(k, \frac{k}{k+R_t}\big)$。
\end{derivationbox}

利用全方差公式，个体继发感染数的方差可直接分解为个体传播泊松随机性与环境伽马随机性两部分：
\begin{align}
  \Var(X_i \mid R_t) &= \E\big[\Var(X_i \mid \nu_i)\big] + \Var\big[\E(X_i \mid \nu_i)\big] \nonumber \\
  &= \E(\nu_i) + \Var(\nu_i) \nonumber \\
  &= R_t + \frac{R_t^2}{k}.
  \label{eq:offspring}
\end{align}
显然，当 $k \to \infty$ 时，$\Var(\nu_i) \to 0$，个体传染力退化为常数 $R_t$，后代分布退化为标准泊松分布 $\text{Poisson}(R_t)$。

\subsection{新发总感染数 $I_t$ 的独立求和性质}

当上一代有 $n$ 名互相独立的感染者时，当期新发总感染数为 $I_t = \sum_{i=1}^n X_{i,t}$。

\begin{derivationbox}{推导 2.2：负二项独立和的分布}
由式~\eqref{eq:nb-pgf}，单一个体的后代 PGF 为：
\begin{equation}
  G_X(s) = \left( \frac{p}{1 - (1-p)s} \right)^k, \qquad p = \frac{k}{k+R_t}.
\end{equation}
由于 $X_{1,t}, \dots, X_{n,t}$ 相互独立，和变量 $I_t$ 的概率母函数为各自母函数之积：
\begin{equation}
  G_I(s) = \E\big(s^{\sum_{i=1}^n X_{i,t}}\big) = \prod_{i=1}^n \E(s^{X_{i,t}}) = [G_X(s)]^n = \left( \frac{p}{1 - (1-p)s} \right)^{nk}.
\end{equation}
这正是参数为分散参数 $nk$、成功概率 $p = \frac{k}{k+R_t}$ 的负二项分布。因此，条件在 $n$ 与 $R_t$ 下：
\begin{equation}
  I_t \mid n, R_t \sim \NB\!\left(nk,\ \frac{k}{k+R_t}\right),
\end{equation}
其理论矩为：
\begin{equation}
  \E(I_t \mid n, R_t) = n R_t, \qquad \Var(I_t \mid n, R_t) = n R_t + \frac{n R_t^2}{k}.
\end{equation}
\end{derivationbox}

\subsection{二项稀疏（Binomial Thinning）下负二项族的闭包定理}

在监测系统只能以概率 $\rho$ 报告病例时，观测病例数 $C_t \mid I_t \sim \Bin(I_t, \rho)$。以下证明：二项稀疏操作严格保持负二项分布族的形式不变，仅平移其参数。

\begin{theorembox}{定理 2.1（二项稀疏下的负二项保持定理）}
若总感染数 $I \sim \NB(r, p)$，且观测数 $C \mid I \sim \Bin(I, \rho)$，其中 $\rho \in (0, 1]$，则边际观测数 $C$ 严格服从负二项分布：
\begin{equation}
  C \sim \NB\left(r,\ \frac{p}{p + \rho(1-p)}\right).
\end{equation}
\end{theorembox}

\begin{proof}
利用重期望法则（Law of iterated expectations）与二项分布的概率母函数。已知对于二项分布 $C \mid I \sim \Bin(I, \rho)$，其条件母函数为：
\begin{equation}
  \E(s^C \mid I) = \sum_{c=0}^I \binom{I}{c} (\rho s)^c (1-\rho)^{I-c} = (1 - \rho + \rho s)^I.
\end{equation}
因此，$C$ 的边际概率母函数 $G_C(s)$ 为：
\begin{equation}
  G_C(s) = \E\big[\E(s^C \mid I)\big] = \E\big[(1 - \rho + \rho s)^I\big] = G_I(1 - \rho + \rho s).
  \label{eq:thin-compound}
\end{equation}
将 $I \sim \NB(r, p)$ 的母函数表达式代入式~\eqref{eq:thin-compound}，将变量 $s$ 整体替换为 $1 - \rho + \rho s$：
\begin{align}
  G_C(s) &= \left( \frac{p}{1 - (1-p)(1 - \rho + \rho s)} \right)^r \nonumber \\
  &= \left( \frac{p}{1 - (1-p)(1-\rho) - (1-p)\rho s} \right)^r.
  \label{eq:thin-expand}
\end{align}
定义辅助分母常数项：
\begin{equation}
  q \equiv 1 - (1-p)(1-\rho) = 1 - (1 - \rho - p + p\rho) = p + \rho(1-p).
\end{equation}
在式~\eqref{eq:thin-expand} 的分母中提出因子 $q$：
\begin{equation}
  G_C(s) = \left( \frac{p}{q - (1-p)\rho s} \right)^r = \left( \frac{p/q}{1 - \frac{(1-p)\rho}{q} s} \right)^r.
  \label{eq:thin-factored}
\end{equation}
注意到：
\begin{equation}
  1 - \frac{p}{q} = \frac{q - p}{q} = \frac{p + \rho(1-p) - p}{q} = \frac{(1-p)\rho}{q}.
\end{equation}
将此关系代入式~\eqref{eq:thin-factored}，得：
\begin{equation}
  G_C(s) = \left( \frac{p'}{1 - (1 - p')s} \right)^r, \qquad \text{其中 } p' \equiv \frac{p}{q} = \frac{p}{p + \rho(1-p)}.
\end{equation}
对照负二项分布概率母函数的定义式~\eqref{eq:nb-pgf}，母函数形式完全相同，分散参数仍为 $r$，成功概率由 $p$ 变换为 $p'$。由于概率母函数与离散概率分布之间存在一一对应性，定理得证。
\end{proof}

\subsection{报告病例数 $C_t$ 的精确条件矩推导}

\begin{derivationbox}{推导 2.3：代入流行病学参数}
在我们的模型中，分散参数为 $r = nk$，原成功概率为 $p = \frac{k}{k + R_t}$，故 $1 - p = \frac{R_t}{k + R_t}$。计算稀疏后的分母常数 $q$ 与新参数 $p'$：
\begin{equation}
  q = p + \rho(1-p) = \frac{k}{k+R_t} + \frac{\rho R_t}{k+R_t} = \frac{k + \rho R_t}{k + R_t},
\end{equation}
\begin{equation}
  p' = \frac{p}{q} = \frac{\frac{k}{k+R_t}}{\frac{k+\rho R_t}{k+R_t}} = \frac{k}{k + \rho R_t}.
\end{equation}
由此得出核心结论：
\begin{equation}
  \boxed{\, C_t \mid n, R_t \sim \NB\left(nk,\ \frac{k}{k + \rho R_t}\right) \,}
  \label{eq:der-ct-dist}
\end{equation}
定义 $m \equiv \rho n$ 为预期报告规模。应用负二项分布矩公式式~\eqref{eq:nb-mean} 与式~\eqref{eq:nb-var}：
\begin{align}
  \E(C_t \mid n, R_t) &= \frac{r(1-p')}{p'} = nk \cdot \frac{\frac{\rho R_t}{k+\rho R_t}}{\frac{k}{k+\rho R_t}} = \rho n R_t = m R_t, \\
  \Var(C_t \mid n, R_t) &= \frac{r(1-p')}{(p')^2} = nk \cdot \frac{\rho R_t}{k+\rho R_t} \cdot \left(\frac{k+\rho R_t}{k}\right)^2 \nonumber \\
  &= \rho n R_t \left(1 + \frac{\rho R_t}{k}\right) = \rho n R_t + \frac{\rho^2 n R_t^2}{k} \nonumber \\
  &= m R_t + \frac{\rho m R_t^2}{k}.
\end{align}
进而得到相对方差（平方变异系数 $\CV^2$）：
\begin{equation}
  \CV^2(C_t \mid n, R_t) = \frac{\Var(C_t \mid n, R_t)}{[\E(C_t \mid n, R_t)]^2} = \frac{\rho n R_t + \rho^2 n R_t^2 / k}{(\rho n R_t)^2} = \frac{1}{\rho n R_t} + \frac{1}{nk} = \frac{1}{m R_t} + \frac{1}{nk}.
  \label{eq:der-cv2}
\end{equation}
\end{derivationbox}

\begin{insightbox}{物理与流行病学洞见 2.1}
式~\eqref{eq:der-cv2} 清晰揭示了观测噪声的内在结构：
\begin{enumerate}
  \item 第一项 $\frac{1}{\rho n R_t} = \frac{1}{m R_t}$ 为\textbf{泊松/二项抽样噪声}。提高报告率 $\rho$ 或增大祖代 $n$ 都能将其压缩；
  \item 第二项 $\frac{1}{nk}$ 为\textbf{超传播个体方差}。请注意：该项分母中\textbf{完全不包含报告率 $\rho$}！这意味着：即使建立 100\% 完美的报告系统（$\rho = 1$），由于群体内超级传播者的存在，每代传播的相对方差依然保留底板 $\frac{1}{nk}$。要消除此项，唯一途径是扩大流行被追踪祖代的绝对人数 $n$。
\end{enumerate}
\end{insightbox}

\section{估计目标 1（条件瞬时再生数 $R_t$）的精确推导}
\label{sec:part03}

\subsection{自然估计量的定义与无偏性}

在给定当期真实环境所确定的瞬时再生数 $R_t$ 时，自然插值估计量定义为观测病例数除以预期报告规模：
\begin{equation}
  \widehat{R}_t \equiv \frac{C_t}{m} = \frac{C_t}{\rho n}.
  \label{eq:estimator-def}
\end{equation}

\begin{derivationbox}{推导 3.1：条件均值与条件方差}
由式~\eqref{eq:der-ct-dist}，条件期望满足：
\begin{equation}
  \E(\widehat{R}_t \mid R_t) = \frac{\E(C_t \mid n, R_t)}{\rho n} = \frac{\rho n R_t}{\rho n} = R_t,
\end{equation}
证明该估计量关于条件目标 $R_t$ 是\textbf{严格无偏}的。

其条件方差为：
\begin{align}
  \Var(\widehat{R}_t \mid R_t) &= \frac{\Var(C_t \mid n, R_t)}{(\rho n)^2} \nonumber \\
  &= \frac{\rho n R_t + \rho^2 n R_t^2 / k}{\rho^2 n^2} \nonumber \\
  &= \frac{R_t}{\rho n} + \frac{R_t^2}{kn} \nonumber \\
  &= \frac{R_t}{m} + \frac{\rho R_t^2}{km} = \frac{R_t (k + \rho R_t)}{k \rho n}.
  \label{eq:var-target1-exact}
\end{align}
\end{derivationbox}

\subsection{对数似然函数与得分函数（Score Function）}

为探究估计量 $\widehat{R}_t$ 的统计效率，我们显式写出给定祖代规模 $n$ 时，观测值 $C_t = c$ 的条件对数似然函数。

\begin{derivationbox}{推导 3.2：得分函数解析分解}
由式~\eqref{eq:der-ct-dist}，观测值 $c$ 的条件概率质量函数为：
\begin{equation}
  P(C_t = c \mid R_t) = \frac{\Gamma(c + nk)}{c! \, \Gamma(nk)} \left( \frac{k}{k + \rho R_t} \right)^{nk} \left( \frac{\rho R_t}{k + \rho R_t} \right)^c.
\end{equation}
取自然对数，得到对数似然函数 $\ell(R_t) \equiv \ln P(C_t = c \mid R_t)$：
\begin{equation}
  \ell(R_t) = \ln\Gamma(c + nk) - \ln(c!) - \ln\Gamma(nk) + nk\ln k + c\ln(\rho R_t) - (c + nk)\ln(k + \rho R_t).
\end{equation}
对参数 $R_t$ 求一阶偏导数，得到得分函数（Score function）$\dot{\ell}(R_t)$：
\begin{align}
  \dot{\ell}(R_t) &= \frac{\partial \ell(R_t)}{\partial R_t} = \frac{c}{R_t} - \frac{(c + nk)\rho}{k + \rho R_t} \nonumber \\
  &= \frac{c(k + \rho R_t) - (c + nk)\rho R_t}{R_t(k + \rho R_t)} \nonumber \\
  &= \frac{ck + c\rho R_t - c\rho R_t - nk\rho R_t}{R_t(k + \rho R_t)} \nonumber \\
  &= \frac{k(c - \rho n R_t)}{R_t(k + \rho R_t)}.
  \label{eq:score-raw}
\end{align}
在式~\eqref{eq:score-raw} 的分子中提取因子 $k\rho n$：
\begin{equation}
  \dot{\ell}(R_t) = \frac{k\rho n}{R_t(k + \rho R_t)} \left( \frac{c}{\rho n} - R_t \right) = \mathcal{I}(R_t) \, (\widehat{R}_t - R_t),
  \label{eq:score-factorized}
\end{equation}
其中定义尺度因子：
\begin{equation}
  \mathcal{I}(R_t) \equiv \frac{k\rho n}{R_t(k + \rho R_t)} = \frac{n}{\frac{R_t}{\rho} + \frac{R_t^2}{k}}.
  \label{eq:fisher-def}
\end{equation}
\end{derivationbox}

\subsection{Fisher 信息量与 Cramér--Rao 下界达到性证明}

\begin{theorembox}{命题 1 的完全证明（Cramér--Rao 下界达到性）}
在假设 H1--H4 下：
\begin{enumerate}[label=(\alph*)]
  \item 观测样本 $C_t$ 关于参数 $R_t$ 的 Fisher 信息量严格等于式~\eqref{eq:fisher-def} 中的 $\mathcal{I}(R_t)$；
  \item 自然估计量 $\widehat{R}_t$ 的条件方差严格达到 Cramér--Rao 下界：
  \begin{equation}
    \Var(\widehat{R}_t \mid R_t) = \frac{1}{\mathcal{I}(R_t)}.
  \end{equation}
\end{enumerate}
\end{theorembox}

\begin{proof}
根据定义，Fisher 信息量是得分函数的二阶原点矩（在满足正则条件下亦等于其方差）：
\begin{align}
  \mathcal{I}_{\text{true}}(R_t) &\equiv \Var\big(\dot{\ell}(R_t) \mid R_t\big) \nonumber \\
  &= \Var\left( \mathcal{I}(R_t)(\widehat{R}_t - R_t) \ \Big|\ R_t \right) \nonumber \\
  &= [\mathcal{I}(R_t)]^2 \, \Var(\widehat{R}_t \mid R_t).
\end{align}
将式~\eqref{eq:var-target1-exact} 的方差代入：
\begin{align}
  \mathcal{I}_{\text{true}}(R_t) &= \left( \frac{k\rho n}{R_t(k + \rho R_t)} \right)^2 \left( \frac{R_t(k + \rho R_t)}{k\rho n} \right) \nonumber \\
  &= \frac{k\rho n}{R_t(k + \rho R_t)} = \mathcal{I}(R_t).
\end{align}
由此立即得出：
\begin{equation}
  \Var(\widehat{R}_t \mid R_t) = \frac{1}{\mathcal{I}(R_t)}.
\end{equation}
根据经典参数统计推断定理 \citep{casella2002statistical}，Cramér--Rao 不等式取等号的\textbf{充要条件}是：对数似然的得分函数可以表示为参数与待估无偏统计量之差的线性函数：
\begin{equation}
  \frac{\partial \ell(R_t)}{\partial R_t} = A(R_t) \big[ T(C_t) - g(R_t) \big].
\end{equation}
对照式~\eqref{eq:score-factorized}，此处 $T(C_t) = \widehat{R}_t = \frac{C_t}{\rho n}$，$g(R_t) = R_t$，$A(R_t) = \mathcal{I}(R_t)$。该条件对\textbf{每一个有限样本量 $n$ 均严格恒等成立}，证明完成。
\end{proof}

\subsection{UMVUE（一致最小方差无偏估计）的唯一性证明}

\begin{theorembox}{定理 3.1（Lehmann--Scheffé 唯一最优性定理）}
在估计目标 1 下，自然估计量 $\widehat{R}_t = C_t/m$ 不仅达到了 Cramér--Rao 界，而且是全参数空间内唯一的均匀最小方差无偏估计（UMVUE）。
\end{theorembox}

\begin{proof}
重写式~\eqref{eq:der-ct-dist} 中 $C_t$ 的概率质量函数：
\begin{align}
  P(C_t = c \mid R_t) &= \frac{\Gamma(c + nk)}{c! \, \Gamma(nk)} \left( \frac{k}{k + \rho R_t} \right)^{nk} \left( \frac{\rho R_t}{k + \rho R_t} \right)^c \nonumber \\
  &= h(c) \exp\left\{ c \ln\left( \frac{\rho R_t}{k + \rho R_t} \right) + nk \ln\left( \frac{k}{k + \rho R_t} \right) \right\},
\end{align}
其中 $h(c) = \frac{\Gamma(c+nk)}{c!\,\Gamma(nk)}$。这表明条件分布族属于标准的\textbf{单参数正则指数族（One-parameter exponential family）}。
其自然参数为：
\begin{equation}
  \upsilon(R_t) \equiv \ln\left( \frac{\rho R_t}{k + \rho R_t} \right) = \ln\left( 1 - \frac{k}{k + \rho R_t} \right).
\end{equation}
当真实再生数 $R_t$ 遍历物理取值区间 $(0, \infty)$ 时，自然参数 $\upsilon$ 单调连续地遍历区间 $(-\infty, 0)$。
因为自然参数空间包含实轴上的非退化开区间，根据指数族的完备性定理 \citep{lehmann2005testing}，充分统计量 $T(C_t) = C_t$ 是\textbf{完备充分统计量（Complete sufficient statistic）}。

由于自然估计量 $\widehat{R}_t = \frac{C_t}{\rho n}$ 是完备充分统计量 $C_t$ 的单变量确定性 Borel 可测函数，且已证 $\E(\widehat{R}_t \mid R_t) = R_t$ 无偏，根据 \textbf{Lehmann--Scheffé 定理}，任何以完备充分统计量为自变量的无偏估计量必为唯一的 UMVUE。证毕。
\end{proof}

\section{估计目标 2（跨期平均再生数 $\bar{R}$）的方差分解}
\label{sec:part04}

\subsection{估计误差的正交分解}

当监测目标从瞬时状态 $R_t$ 转向跨期总体均值 $\Rbar \equiv \E(R_t)$ 时，估计误差的基准参照系发生根本转变。自然估计量 $\widehat{R}_t = C_t/m$ 的总误差可严格拆分为两部分：
\begin{equation}
  \widehat{R}_t - \Rbar = \underbrace{(\widehat{R}_t - R_t)}_{\text{抽样与观测误差}} + \underbrace{(R_t - \Rbar)}_{\text{环境外生波动误差}}.
  \label{eq:error-split}
\end{equation}

\begin{derivationbox}{推导 4.1：正交交叉项为零的严格证明}
计算总误差的均方值（MSE）：
\begin{align}
  \E\big[(\widehat{R}_t - \Rbar)^2 \mid n\big] &= \E\Big\{ \big[(\widehat{R}_t - R_t) + (R_t - \Rbar)\big]^2 \ \Big|\ n \Big\} \nonumber \\
  &= \E\big[(\widehat{R}_t - R_t)^2 \mid n\big] + \E\big[(R_t - \Rbar)^2 \mid n\big] + 2\,\E\big[(\widehat{R}_t - R_t)(R_t - \Rbar) \mid n\big].
\end{align}
考察交叉项。利用重期望法则，先对给定当期环境 $R_t$ 取条件期望：
\begin{align}
  \E\big[(\widehat{R}_t - R_t)(R_t - \Rbar) \mid n\big] &= \E\Big\{ \E\big[(\widehat{R}_t - R_t)(R_t - \Rbar) \mid n, R_t\big] \ \Big|\ n \Big\} \nonumber \\
  &= \E\Big\{ (R_t - \Rbar) \, \E(\widehat{R}_t - R_t \mid n, R_t) \ \Big|\ n \Big\}.
\end{align}
根据推导 3.1 中证明的无偏性，$\E(\widehat{R}_t - R_t \mid n, R_t) = 0$ 恒成立。因此：
\begin{equation}
  \E\big[(\widehat{R}_t - R_t)(R_t - \Rbar) \mid n\big] = \E\big[ (R_t - \Rbar) \cdot 0 \mid n \big] \equiv 0.
\end{equation}
这证明抽样误差与环境波动在概率空间中\textbf{严格正交}，二者协方差恒等于零。因此总均方误差等于方差之和：
\begin{equation}
  \E\big[(\widehat{R}_t - \Rbar)^2 \mid n\big] = \E\big[(\widehat{R}_t - R_t)^2 \mid n\big] + \Var(R_t \mid n).
\end{equation}
\end{derivationbox}

\subsection{全方差公式展开与闭式解 $v_R + A/m$}

\begin{theorembox}{命题 2 的完全证明（跨期平均估计的方差分解）}
在假设 H1--H4 下，自然估计量 $\widehat{R}_t = C_t / (\rho n)$ 关于跨期均值 $\Rbar$ 的边际方差服从解析闭式：
\begin{equation}
  V(m) \equiv \Var(\widehat{R}_t \mid n) = v_R + \frac{A}{m},
  \label{eq:var-closed}
\end{equation}
其中环境方差 $v_R = \Var(R_t)$，预期报告规模 $m = \rho n$，规模项系数 $A$ 严格由下式给出：
\begin{equation}
  A \equiv \Rbar + \frac{\rho(\Rbar^2 + v_R)}{k}.
  \label{eq:A-closed}
\end{equation}
\end{theorembox}

\begin{proof}
根据全方差公式（Eve's law）：
\begin{equation}
  \Var(\widehat{R}_t \mid n) = \underbrace{\E\big[\Var(\widehat{R}_t \mid n, R_t) \mid n\big]}_{\text{条件方差的期望（抽样项）}} + \underbrace{\Var\big[\E(\widehat{R}_t \mid n, R_t) \mid n\big]}_{\text{条件均值的方差（环境项）}}.
  \label{eq:eve-law}
\end{equation}
我们分别计算这两项：
\begin{enumerate}
  \item \textbf{第二项（环境方差项）：}
  由条件无偏性，$\E(\widehat{R}_t \mid n, R_t) = R_t$。结合假设 H4（祖代规模 $n$ 外生独立于环境扰动）：
  \begin{equation}
    \Var\big[\E(\widehat{R}_t \mid n, R_t) \mid n\big] = \Var(R_t \mid n) = \Var(R_t) \equiv v_R.
  \end{equation}
  此项不包含任何样本规模变量（既无 $n$，亦无 $\rho$），构成了方差下确界。

  \item \textbf{第一项（抽样随机性在环境上的边际平均）：}
  将式~\eqref{eq:var-target1-exact} 代入外层期望：
  \begin{align}
    \E\big[\Var(\widehat{R}_t \mid n, R_t) \mid n\big] &= \E\left[ \frac{R_t}{\rho n} + \frac{R_t^2}{kn} \ \Bigg|\ n \right] \nonumber \\
    &= \frac{\E(R_t \mid n)}{\rho n} + \frac{\E(R_t^2 \mid n)}{kn}.
    \label{eq:step-sampling-exp}
  \end{align}
  利用方差与二阶原点矩的基本关系：
  \begin{equation}
    \E(R_t^2 \mid n) = \Var(R_t \mid n) + [\E(R_t \mid n)]^2 = v_R + \Rbar^2.
  \end{equation}
  将其代入式~\eqref{eq:step-sampling-exp}：
  \begin{equation}
    \E\big[\Var(\widehat{R}_t \mid n, R_t) \mid n\big] = \frac{\Rbar}{\rho n} + \frac{\Rbar^2 + v_R}{kn}.
  \end{equation}
  提取公分母 $m = \rho n$：
  \begin{align}
    \frac{\Rbar}{\rho n} + \frac{\Rbar^2 + v_R}{kn} &= \frac{\Rbar}{m} + \frac{\rho(\Rbar^2 + v_R)}{k(\rho n)} \nonumber \\
    &= \frac{1}{m} \left[ \Rbar + \frac{\rho(\Rbar^2 + v_R)}{k} \right] \equiv \frac{A}{m}.
  \end{align}
\end{enumerate}
将两项相加，即得 $V(m) = v_R + \frac{A}{m}$。
\end{proof}

\begin{insightbox}{数学细节洞见 4.1：为什么超传播项中含有 $v_R$？}
仔细观察 $A$ 的表达式：$A = \Rbar + \frac{\rho(\Rbar^2 + v_R)}{k}$。
很多初学者容易误以为分子应当是 $\Rbar^2$。为什么环境方差 $v_R$ 会再次出现在抽样项分子中？
其数学根源在于：负二项分布的方差正比于均值的\textbf{二次方}（$R_t^2/k$）。
由于平方函数是凸函数，根据 Jensen 不等式，环境扰动使得 $\E(R_t^2) = \Rbar^2 + v_R > \Rbar^2$。
当暴发环境出现极端高值（如 $R_t \gg 1$）时，超传播引发的病例聚集呈现非线性的爆炸式放大。因此，环境方差本身不仅作为常数平移量 $v_R$ 进入底板，还通过过离散机制\textbf{反向放大了抽样项本身的离散度}！
\end{insightbox}

\subsection{对 $R_t$ 条件无偏估计类的通用最优性界限}

\begin{theorembox}{定理 4.1（条件无偏估计类的统一最小均方误差下界）}
设 $\widetilde{\Theta}$ 为满足 $\E(\tilde{\theta} \mid R_t) = R_t$ 的所有单期估计量构成的集合。则对任意 $\tilde{\theta} \in \widetilde{\Theta}$，其估计跨期均值 $\Rbar$ 的均方误差必然满足：
\begin{equation}
  \E\big[(\tilde{\theta} - \Rbar)^2\big] \ge v_R + \frac{A}{m},
\end{equation}
且当且仅当 $\tilde{\theta} = \widehat{R}_t = C_t/m$ 时取等号。
\end{theorembox}

\begin{proof}
根据推导 4.1 的正交分解，对于任意 $\tilde{\theta} \in \widetilde{\Theta}$：
\begin{equation}
  \E\big[(\tilde{\theta} - \Rbar)^2\big] = \Var(R_t) + \E\big[(\tilde{\theta} - R_t)^2\big] = v_R + \E\big[ \Var(\tilde{\theta} \mid R_t) \big].
\end{equation}
根据命题 1，给定 $R_t$ 时，由单参数指数族 Cramér--Rao 定理，任何条件无偏估计量的条件方差均受 Fisher 信息倒数的下界约束：
\begin{equation}
  \Var(\tilde{\theta} \mid R_t) \ge \frac{1}{\mathcal{I}(R_t)} = \frac{R_t}{\rho n} + \frac{R_t^2}{kn}.
\end{equation}
两侧对环境分布取数学期望：
\begin{equation}
  \E\big[ \Var(\tilde{\theta} \mid R_t) \big] \ge \E\left[ \frac{1}{\mathcal{I}(R_t)} \right] = \E\left[ \frac{R_t}{\rho n} + \frac{R_t^2}{kn} \right] = \frac{A}{m}.
\end{equation}
两端加上常数 $v_R$，即得出下界：
\begin{equation}
  \E\big[(\tilde{\theta} - \Rbar)^2\big] \ge v_R + \frac{A}{m}.
\end{equation}
由于自然估计量 $\widehat{R}_t$ 处处达到此条件下界，等号成立。证毕。
\end{proof}

\subsection{一步前瞻预测的最优性证明}

\begin{derivationbox}{推导 4.2：预测下界式~\eqref{eq:predict-bound} 的展开}
考虑预测问题：根据截至第 $t$ 期的历史信息集 $\mathcal{F}_t$，预测下一代经标准化的新发报告量 $y = C_{t+1}/m$。
设 $\hat{y} \in \mathcal{F}_t$ 为任意前瞻预测量。在假设 H4 下，$R_{t+1}$ 独立于 $\mathcal{F}_t$。
利用预测均方误差的标准偏置-方差分解：
\begin{align}
  \E\big[(\hat{y} - y)^2 \mid \mathcal{F}_t\big] &= \E\big[(\hat{y} - \E(y \mid \mathcal{F}_t) + \E(y \mid \mathcal{F}_t) - y)^2 \mid \mathcal{F}_t\big] \nonumber \\
  &= \big(\hat{y} - \E(y \mid \mathcal{F}_t)\big)^2 + \Var(y \mid \mathcal{F}_t).
\end{align}
由独立性，$\E(y \mid \mathcal{F}_t) = \E(C_{t+1}/m) = \Rbar$。因此：
\begin{equation}
  \Var(y \mid \mathcal{F}_t) = \Var\left( \frac{C_{t+1}}{m} \right) = V(m) = v_R + \frac{A}{m}.
\end{equation}
故预测误差为：
\begin{equation}
  \E\big[(\hat{y} - y)^2 \mid \mathcal{F}_t\big] = (\hat{y} - \Rbar)^2 + v_R + \frac{A}{m} \ge v_R + \frac{A}{m}.
  \label{eq:predict-bound}
\end{equation}
当且仅当选择前瞻预测量 $\hat{y} = \Rbar$ 时，第一项非负平方项被压缩至 0，达到理论预测极限。
\end{derivationbox}

\section{方差平衡规模 $m_\times$、边际收益与非对称性}
\label{sec:part05}

\subsection{方差平衡规模的定义与物理几何意义}

在方差闭式解 $V(m) = v_R + A/m$ 中，可约的抽样随机性 $A/m$ 随规模线性衰减，而不可约的环境方差 $v_R$ 保持恒定。定义二者相等的交点规模：

\begin{definition}[方差平衡规模 $\mx$]
定义使得可约抽样方差与不可约环境方差严格相等的特征观测规模为：
\begin{equation}
  \frac{A}{\mx} = v_R \quad \iff \quad \mx \equiv \frac{A}{v_R} = \frac{\Rbar + \frac{\rho(\Rbar^2 + v_R)}{k}}{v_R}.
  \label{eq:mx-def}
\end{equation}
\end{definition}

借助平衡规模 $\mx$，总方差可紧凑表征为无量纲规模比值 $m / \mx$ 的函数：
\begin{equation}
  V(m) = v_R \left( 1 + \frac{\mx}{m} \right).
  \label{eq:var-compact}
\end{equation}

\begin{insightbox}{审慎定位：平衡规模不是临床或预警的“充分阈值”}
本手册与论文在此做出严格强调：方差平衡规模 $\mx$ 仅表示抽样项与环境项相等的方差结构转折基准（Balance point），\textbf{绝非}满足特定估计精度或预警检测功效要求的最低实用规模（Minimum operational threshold）。在不同公共卫生任务中：
\begin{itemize}
  \item 若目标为以 80\% 功效检测 15\% 的传播增长，所需规模 $m \approx 527$，约为 $\mx$ 的 10 倍（见第~\ref{sec:part11} 节）；
  \item 若目标为压制跨期宏观趋势估计误差，当规模达到 $\mx$ 时，继续扩样的边际收益已大幅钝化，资源应当向增加期数或空间池化转移。
\end{itemize}
\end{insightbox}

\subsection{边际收益递减的四大数学定理证明}

\begin{derivationbox}{推导 5.1：方差份额、翻倍收益与弹性分析}
基于式~\eqref{eq:var-compact}，我们可以严格导出方差曲线的四个关键几何定理：

\begin{enumerate}
  \item \textbf{抽样方差占比定理：}
  抽样随机性在总方差中所占的百分比为：
  \begin{equation}
    S(m) \equiv \frac{A/m}{V(m)} = \frac{A/m}{v_R + A/m} = \frac{\frac{A/v_R}{m}}{1 + \frac{A/v_R}{m}} = \frac{\mx/m}{1 + \mx/m} = \frac{1}{1 + \frac{m}{\mx}}.
    \label{eq:share-proof}
  \end{equation}
  当 $m = \mx$ 时，$S(\mx) = \frac{1}{1 + 1} = \frac{1}{2} = 50\%$。两类方差严格均分总误差。

  \item \textbf{规模翻倍收益衰减定理（Doubling reduction）：}
  将单期监测规模从 $m$ 翻倍扩大到 $2m$ 时，总方差相对下降的比例为：
  \begin{align}
    \Delta_2(m) &\equiv \frac{V(m) - V(2m)}{V(m)} = \frac{\left(v_R + \frac{A}{m}\right) - \left(v_R + \frac{A}{2m}\right)}{v_R + \frac{A}{m}} \nonumber \\
    &= \frac{\frac{A}{2m}}{v_R + \frac{A}{m}} = \frac{1}{2} \cdot \frac{A/m}{v_R + A/m} = \frac{1}{2} S(m) = \frac{1}{2\left(1 + \frac{m}{\mx}\right)}.
    \label{eq:doubling-proof}
  \end{align}
  考察三个典型区域：
  \begin{itemize}
    \item 当 $m \to 0$（小样本极端）时，$\Delta_2 \to 50\%$，符合经典泊松抽样方差减半直觉；
    \item 当 $m = \mx$（平衡点）时，$\Delta_2 = \frac{1}{4} = 25.0\%$；
    \item 当 $m = 4\mx$（大样本饱和）时，$\Delta_2 = \frac{1}{2(1+4)} = 10.0\%$，翻倍投入仅换来不到一成的方差压缩。
  \end{itemize}

  \item \textbf{规模弹性（Elasticity）定理：}
  总方差关于监测规模 $m$ 的点弹性为：
  \begin{equation}
    \epsilon_m \equiv \frac{\mathrm{d}\ln V(m)}{\mathrm{d}\ln m} = \frac{m}{V(m)} \frac{\mathrm{d}V(m)}{\mathrm{d}m} = \frac{m}{v_R + A/m} \left( -\frac{A}{m^2} \right) = -\frac{A/m}{v_R + A/m} = -\frac{\mx}{m + \mx}.
  \end{equation}
  在平衡点 $m = \mx$ 处，规模弹性严格等于 $-\frac{1}{2}$。

  \item \textbf{处处严格凸性（Strict convexity）：}
  求 $V(m)$ 关于 $m$ 的一阶与二阶导数：
  \begin{equation}
    V'(m) = -\frac{A}{m^2} < 0, \qquad V''(m) = \frac{2A}{m^3} > 0, \quad \forall m > 0.
  \end{equation}
  一阶导数处处严格为负，二阶导数在正实轴上处处严格为正。函数没有拐点，曲线呈现处处向原点凸起的双曲线单调下降特征。
\end{enumerate}
\end{derivationbox}

\subsection{覆盖面与报告率的非对称边际效应（Remark 1 证明）}

预期报告规模由两个途径共同构成：$m = \rho n$。究竟应该追求让既定人群的报告更全面（提高报告概率 $\rho$），还是覆盖更多的人口群体（扩大祖代样本 $n$）？

\begin{theorembox}{注 1 的完全证明（覆盖面与报告率的边际非对称性）}
在假设 H1--H4 下：
\begin{enumerate}[label=(\roman*)]
  \item 若固定祖代规模 $n$ 而令报告率 $\rho \to 1$，方差趋于大于 $v_R$ 的严格常数；
  \item 若固定报告率 $\rho$ 而令覆盖规模 $n \to \infty$，方差趋于理论底板 $v_R$；
  \item 按相同百分比增加 $n$ 所取得的方差降幅，恒等于增加 $\rho$ 所取得降幅的 $A/\Rbar$ 倍：
  \begin{equation}
    \frac{-\frac{\partial V}{\partial \ln n}}{-\frac{\partial V}{\partial \ln \rho}} = \frac{A}{\Rbar} = 1 + \frac{\rho(\Rbar^2 + v_R)}{k\Rbar} > 1.
    \label{eq:asym-ratio}
  \end{equation}
\end{enumerate}
\end{theorembox}

\begin{proof}
将方差式~\eqref{eq:var-closed} 展开为 $n$ 和 $\rho$ 的显式二元函数：
\begin{equation}
  V(n, \rho) = v_R + \frac{\Rbar}{\rho n} + \frac{\Rbar^2 + v_R}{kn}.
\end{equation}
\begin{enumerate}
  \item 固定 $n$ 令 $\rho \to 1$：
  \begin{equation}
    \lim_{\rho \to 1} V(n, \rho) = v_R + \frac{\Rbar}{n} + \frac{\Rbar^2 + v_R}{kn} = v_R + \frac{1}{n}\left(\Rbar + \frac{\Rbar^2+v_R}{k}\right) > v_R.
  \end{equation}
  因超传播项不随 $\rho$ 衰减，存在由有限祖代规模 $n$ 锁定的残余人口随机性。
  \item 固定 $\rho$ 令 $n \to \infty$：
  \begin{equation}
    \lim_{n \to \infty} V(n, \rho) = v_R + 0 + 0 = v_R.
  \end{equation}
  \item 计算对数偏导数（即方差对比例变化的敏感度）：
  \begin{align}
    -\frac{\partial V}{\partial \ln n} &= -n \frac{\partial V}{\partial n} = -n \left[ -\frac{\Rbar}{\rho n^2} - \frac{\Rbar^2 + v_R}{kn^2} \right] = \frac{\Rbar}{\rho n} + \frac{\Rbar^2 + v_R}{kn} = \frac{A}{m}, \\
    -\frac{\partial V}{\partial \ln \rho} &= -\rho \frac{\partial V}{\partial \rho} = -\rho \left[ -\frac{\Rbar}{\rho^2 n} \right] = \frac{\Rbar}{\rho n} = \frac{\Rbar}{m}.
  \end{align}
  两者之比为：
  \begin{equation}
    \frac{-\partial V / \partial \ln n}{-\partial V / \partial \ln \rho} = \frac{A/m}{\Rbar/m} = \frac{A}{\Rbar} = \frac{\Rbar + \frac{\rho(\Rbar^2+v_R)}{k}}{\Rbar} = 1 + \frac{\rho(\Rbar^2 + v_R)}{k\Rbar}.
  \end{equation}
  由于 $\rho > 0, k > 0, \Rbar > 0, v_R \ge 0$，第二项严格为正，故比值恒大于 1。证毕。
\end{enumerate}
\end{proof}

\section{报告异质性、三项方差分解与不可识别性}
\label{sec:part06}

\subsection{随机报告率与两级全方差展开}

在实际疾病监测中，就诊意愿、检测准入及机构填报口径往往随时间剧烈波动。假设 H5 引入随机报告概率 $p_t \in (0, 1]$，其均值为 $\E(p_t) = \rho$，方差为 $\Var(p_t) = \sigma^2$。相对波动率定义为变异系数 $\tau \equiv \sigma / \rho$。

\begin{theorembox}{命题 3 的完全推导（报告异质性下的方差三分量定理）}
在假设 H2--H6 下，自然估计量 $\widehat{R}_t = C_t / (\rho n)$ 的边际方差由三项构成：
\begin{equation}
  V_H(m) = \underbrace{v_R}_{\text{环境方差}} + \underbrace{v_H}_{\text{报告异质性方差}} + \underbrace{\frac{A_H}{m}}_{\text{规模抽样项}},
  \label{eq:vH-closed}
\end{equation}
其中报告异质性方差 $v_H$ 与修正尺度系数 $A_H$ 分别为：
\begin{align}
  v_H &\equiv \tau^2 (\Rbar^2 + v_R), \label{eq:vH-def} \\
  A_H &\equiv \Rbar + \frac{\rho(1+\tau^2)(\Rbar^2 + v_R)}{k} = \Rbar + \frac{\rho U}{k}, \label{eq:AH-def}
\end{align}
式中定义复合参数 $U \equiv (1+\tau^2)(\Rbar^2 + v_R)$。
\end{theorembox}

\begin{proof}
推导需分两层嵌套使用全方差公式：第一层对随机报告率 $p_t$ 条件化，第二层对环境随机性 $R_t$ 条件化。

\paragraph{第一层：给定 $(n, R_t)$ 时对 $p_t$ 展开。}
由式~\eqref{eq:der-ct-dist}，在给定当前报告率 $p_t$ 的微观实现值时，报告数服从 $\NB(nk, \frac{k}{k+p_t R_t})$。条件矩为：
\begin{equation}
  \E(C_t \mid n, R_t, p_t) = p_t n R_t, \qquad \Var(C_t \mid n, R_t, p_t) = p_t n R_t + \frac{p_t^2 n R_t^2}{k}.
\end{equation}
应用全方差公式消除随机变量 $p_t$：
\begin{align}
  \Var(C_t \mid n, R_t) &= \E\big[ \Var(C_t \mid n, R_t, p_t) \mid n, R_t \big] + \Var\big[ \E(C_t \mid n, R_t, p_t) \mid n, R_t \big] \nonumber \\
  &= \E\left[ p_t n R_t + \frac{p_t^2 n R_t^2}{k} \ \Bigg|\ n, R_t \right] + \Var(p_t n R_t \mid n, R_t) \nonumber \\
  &= n R_t \, \E(p_t) + \frac{n R_t^2}{k} \, \E(p_t^2) + n^2 R_t^2 \, \Var(p_t).
\end{align}
代入统计矩 $\E(p_t) = \rho$，$\Var(p_t) = \sigma^2$，以及 $\E(p_t^2) = \Var(p_t) + [\E(p_t)]^2 = \sigma^2 + \rho^2$：
\begin{equation}
  \Var(C_t \mid n, R_t) = \rho n R_t + \frac{(\rho^2 + \sigma^2) n R_t^2}{k} + \sigma^2 n^2 R_t^2.
  \label{eq:var-step1}
\end{equation}
注意式~\eqref{eq:var-step1} 最后一项包含 $n^2$，与祖代规模的平方成正比。

\paragraph{第二层：消除环境随机性 $R_t$。}
考察估计量 $\widehat{R}_t = \frac{C_t}{\rho n}$。其关于 $R_t$ 的条件期望为：
\begin{equation}
  \E(\widehat{R}_t \mid n, R_t) = \frac{\E(C_t \mid n, R_t)}{\rho n} = \frac{\rho n R_t}{\rho n} = R_t.
\end{equation}
再次应用全方差公式：
\begin{align}
  \Var(\widehat{R}_t \mid n) &= \E\big[ \Var(\widehat{R}_t \mid n, R_t) \mid n \big] + \Var\big[ \E(\widehat{R}_t \mid n, R_t) \mid n \big] \nonumber \\
  &= \frac{1}{\rho^2 n^2} \, \E\big[ \Var(C_t \mid n, R_t) \mid n \big] + \Var(R_t \mid n).
\end{align}
将式~\eqref{eq:var-step1} 代入期望项，并使用 $\E(R_t) = \Rbar$ 与 $\E(R_t^2) = \Rbar^2 + v_R$：
\begin{align}
  \frac{1}{\rho^2 n^2} \E\big[ \Var(C_t \mid n, R_t) \mid n \big] &= \frac{1}{\rho^2 n^2} \left[ \rho n \Rbar + \frac{(\rho^2 + \sigma^2) n (\Rbar^2 + v_R)}{k} + \sigma^2 n^2 (\Rbar^2 + v_R) \right] \nonumber \\
  &= \frac{\Rbar}{\rho n} + \frac{(\rho^2 + \sigma^2)(\Rbar^2 + v_R)}{k \rho^2 n} + \frac{\sigma^2}{\rho^2}(\Rbar^2 + v_R).
  \label{eq:var-step2}
\end{align}
利用变异系数定义 $\tau = \sigma / \rho$，式~\eqref{eq:var-step2} 第三项化为：
\begin{equation}
  \frac{\sigma^2}{\rho^2}(\Rbar^2 + v_R) = \tau^2(\Rbar^2 + v_R) \equiv v_H.
\end{equation}
注意此项中的 $n^2$ 与分母 $\rho^2 n^2$ 完全约掉，因此它是一个\textbf{与样本量 $n$ 和期望规模 $m$ 彻底无关的恒定方差项}！

合并前两项并提取分母 $m = \rho n$：
\begin{align}
  \frac{\Rbar}{\rho n} + \frac{(\rho^2 + \sigma^2)(\Rbar^2 + v_R)}{k \rho^2 n} &= \frac{1}{\rho n} \left[ \Rbar + \frac{\rho^2(1 + \tau^2)(\Rbar^2 + v_R)}{k \rho} \right] \nonumber \\
  &= \frac{1}{m} \left[ \Rbar + \frac{\rho(1 + \tau^2)(\Rbar^2 + v_R)}{k} \right] \equiv \frac{A_H}{m}.
\end{align}
加上条件均值的方差 $\Var(R_t \mid n) = v_R$，最终得到：
\begin{equation}
  V_H(m) = v_R + v_H + \frac{A_H}{m}.
\end{equation}
证毕。
\end{proof}

\subsection{推广的方差平衡规模 $m_{\times, H}$}

由于报告异质性在零刻度引入了额外的误差底板 $v_H$，不可约总方差抬升为 $v_R + v_H$。令可约抽样项与此不可约总底板相等：
\begin{equation}
  \frac{A_H}{\mxH} = v_R + v_H \quad \iff \quad \mxH \equiv \frac{A_H}{v_R + v_H}.
\end{equation}

\begin{derivationbox}{推导 6.1：报告异质性压缩平衡规模的证明}
我们证明：在任意非平凡报告波动（$\tau > 0$）下，修正后的平衡规模严格小于无报告波动的理想平衡规模，即 $\mxH < \mx$。
作商计算：
\begin{align}
  \frac{\mxH}{\mx} &= \frac{A_H / (v_R + v_H)}{A / v_R} = \frac{A_H}{A} \cdot \frac{v_R}{v_R + v_H} \nonumber \\
  &= \frac{\Rbar + \frac{\rho(1+\tau^2)(\Rbar^2+v_R)}{k}}{\Rbar + \frac{\rho(\Rbar^2+v_R)}{k}} \cdot \frac{v_R}{v_R + \tau^2(\Rbar^2+v_R)}.
\end{align}
记常数 $C_1 = \Rbar > 0, C_2 = \frac{\rho(\Rbar^2+v_R)}{k} > 0, C_3 = \Rbar^2+v_R > 0$。上式可写为：
\begin{equation}
  \frac{\mxH}{\mx} = \frac{C_1 + C_2(1+\tau^2)}{C_1 + C_2} \cdot \frac{v_R}{v_R + C_3 \tau^2} = \left( 1 + \frac{C_2 \tau^2}{C_1 + C_2} \right) \left( \frac{1}{1 + \frac{C_3}{v_R} \tau^2} \right).
\end{equation}
比较分子相对增长率与分母相对增长率：
\begin{equation}
  \frac{C_2}{C_1 + C_2} = \frac{\frac{\rho(\Rbar^2+v_R)}{k}}{\Rbar + \frac{\rho(\Rbar^2+v_R)}{k}} < 1, \qquad \text{而} \quad \frac{C_3}{v_R} = \frac{\Rbar^2+v_R}{v_R} = 1 + \frac{\Rbar^2}{v_R} > 1.
\end{equation}
由于 $\frac{C_2}{C_1+C_2} < 1 < \frac{C_3}{v_R}$，分母随 $\tau^2$ 膨胀的速度远快于分子。因此对任意 $\tau > 0$：
\begin{equation}
  \mxH < \mx.
\end{equation}
在典型参数下（$\Rbar=1.15, v_R=0.04, k=0.35, \rho=0.25, \tau=0.15$），$\mx$ 从 $53.1$ 骤降至 $\mxH = 30.4$。
\end{derivationbox}

\begin{insightbox}{物理洞见 6.1：报告波动让采样规模更快进入“收益饱和”}
这一结论具有极其深刻的政策含义：\textbf{数据越脏（报告异质性 $\tau$ 越大），扩大病例规模 $m$ 的有效空间就越小}！
因为报告噪声本身构成了坚硬的壁垒，当规模仅达到 30 例左右时，抽样误差就已被系统性报告波动所掩盖。在此之后继续砸钱追踪上报更多病例，是在用昂贵的边际成本去稀释已经微不足道的微观抽样项。
\end{insightbox}

\subsection{统计不可识别性与方差底板的界定}

\begin{theorembox}{命题 4 的证明（边际矩意义下的不可识别性）}
在假设 H1--H6 下，若仅依据单期病例计数的边际均值与方差，环境方差 $v_R$ 与报告变异系数 $\tau^2$ \textbf{不能分别识别}；两类波动在单期边际二阶矩中完全简并为单一复合参数 $U$。
\end{theorembox}

\begin{proof}
根据式~\eqref{eq:vH-closed}、式~\eqref{eq:vH-def} 与式~\eqref{eq:AH-def}，考察单期计数的一阶矩与二阶矩依赖的参数：
\begin{enumerate}
  \item 边际期望：$\E(C_t / m) = \Rbar$。此项仅包含均值 $\Rbar$，不包含方差参数 $v_R$ 与 $\tau^2$。
  \item 边际方差：
  \begin{equation}
    V_H(m) = (v_R + v_H) + \frac{A_H}{m} = \big[ (1+\tau^2)(\Rbar^2+v_R) - \Rbar^2 \big] + \frac{1}{m}\left[ \Rbar + \frac{\rho(1+\tau^2)(\Rbar^2+v_R)}{k} \right].
  \end{equation}
  定义复合标量参数：
  \begin{equation}
    U \equiv (1+\tau^2)(\Rbar^2 + v_R).
  \end{equation}
  则总底板为 $v_R + v_H = U - \Rbar^2$，规模项分子为 $A_H = \Rbar + \frac{\rho U}{k}$。二者均\textbf{仅仅依赖于标量 $U$}。
\end{enumerate}
因此，只要保持复合值 $U = (1+\tau^2)(\Rbar^2+v_R)$ 恒定，参数对 $(v_R, \tau^2)$ 可以在双曲线等值线上任意滑动，而生成的观测计数的单期均值与方差完全相同。故二者在单期截面矩意义下不能分别识别。

时间序列的动态相关结构能否帮助分离两类波动，取决于对环境过程与报告过程施加的额外结构性假设（如第~\ref{sec:part08} 节的分母滞后权重结构或第~\ref{sec:part09} 节的自回归动力学形式）。
\end{proof}

\begin{derivationbox}{推导 6.2：自然估计量在报告异质性下的渐近方差底板}
对于自然估计量 $\widehat{R}_t = C_t / (\rho n)$，当固定报告波动过程并令单期祖代规模 $n \to \infty$（从而预期报告规模 $m = \rho n \to \infty$）时：
\begin{equation}
  \frac{A_H}{m} = \frac{\Rbar + \frac{\rho(1+\tau^2)(\Rbar^2+v_R)}{k}}{m} \longrightarrow 0.
\end{equation}
因此：
\begin{equation}
  \lim_{m \to \infty} V_H(m) = v_R + v_H = v_R + \tau^2(\Rbar^2 + v_R).
\end{equation}
这表明：在存在跨期报告率异质性时，自然估计量 $\widehat{R}_t$ 的均方误差无法随单期规模 $m$ 的扩大而收敛至 0，其保留的渐近方差底板严格为 $v_R + v_H$。

我们指出：该极限是针对特定估计量 $\widehat{R}_t = C_t/m$ 得到的方差下限，反映了直接使用未校正报告计数时单期扩样的局限性；它并不自动构成全参数空间内所有可能估计方法的通用信息论下界。若能借助外部辅助数据（如各辖区填报完整性比例）对报告率波动进行显式校正，则可进一步剥离报告异质性方差 $v_H$。
\end{derivationbox}

\section{决策反转定理：提高报告率与扩大覆盖面的排序}
\label{sec:part07}

\subsection{边际收益差值的显式导数展开}

在考虑报告异质性 $V_H$ 时，公共卫生决策者再次面临路线选择：是将有限资金用于激励医疗机构提高填报完整度 $\rho$，还是扩大监测覆盖的哨点网络与人群 $n$？

\begin{theorembox}{推论 1 的完全证明（策略优劣反转定理）}
将报告相对方差 $\tau^2$ 视为平均报告率 $\rho$ 的泛函 $\tau^2(\rho)$。在假设 H2--H6 下：
\begin{enumerate}[label=(\roman*)]
  \item 提高报告率相较于扩大覆盖面的相对边际方差压制优势严格取决于判别式：
  \begin{equation}
    \Delta_{\rho, n} \equiv -\rho\frac{\partial V_H}{\partial \rho} - \left(-n\frac{\partial V_H}{\partial n}\right) = \frac{\Rbar^2 + v_R}{kn} \left[ -\rho\frac{\mathrm{d}\tau^2}{\mathrm{d}\rho}(kn + 1) - (1 + \tau^2) \right].
    \label{eq:reversal-condition}
  \end{equation}
  \item 若报告波动源于制度性或政策性切换（$\frac{\mathrm{d}\tau^2}{\mathrm{d}\rho} = 0$），扩大覆盖面 $n$ \textbf{恒处严格绝对占优地位}。
  \item 若报告波动源于二项/贝塔抽样抽样误差（$\tau^2 = \frac{1-\rho}{\rho(\nu+1)}$），当且仅当预期规模超过临界阈值时发生反转：
  \begin{equation}
    \boxed{\, m > \frac{\rho^2 \nu}{k} \,}
    \label{eq:reversal-threshold}
  \end{equation}
  在临界点以上，提高报告率 $\rho$ 的边际价值反超扩大覆盖面 $n$。
\end{enumerate}
\end{theorembox}

\begin{proof}
将方差公式式~\eqref{eq:vH-closed} 显式展开为 $n$ 与 $\rho$ 的函数：
\begin{equation}
  V_H(n, \rho) = v_R + \tau^2(\rho)(\Rbar^2 + v_R) + \frac{\Rbar}{\rho n} + \frac{[1 + \tau^2(\rho)](\Rbar^2 + v_R)}{kn}.
\end{equation}
分别计算对数偏导数（即对比例扩增的边际弹性贡献）：
\begin{align}
  -\rho\frac{\partial V_H}{\partial \rho} &= -\rho \left[ \frac{\mathrm{d}\tau^2}{\mathrm{d}\rho}(\Rbar^2+v_R) - \frac{\Rbar}{\rho^2 n} + \frac{\frac{\mathrm{d}\tau^2}{\mathrm{d}\rho}(\Rbar^2+v_R)}{kn} \right] \nonumber \\
  &= -\rho\frac{\mathrm{d}\tau^2}{\mathrm{d}\rho}(\Rbar^2+v_R) + \frac{\Rbar}{\rho n} - \frac{\rho\frac{\mathrm{d}\tau^2}{\mathrm{d}\rho}(\Rbar^2+v_R)}{kn}, \\
  -n\frac{\partial V_H}{\partial n} &= -n \left[ -\frac{\Rbar}{\rho n^2} - \frac{(1+\tau^2)(\Rbar^2+v_R)}{kn^2} \right] \nonumber \\
  &= \frac{\Rbar}{\rho n} + \frac{(1+\tau^2)(\Rbar^2+v_R)}{kn}.
\end{align}
计算边际效益差 $\Delta_{\rho, n} = -\rho\frac{\partial V_H}{\partial \rho} - (-n\frac{\partial V_H}{\partial n})$。
注意到公共项 $\frac{\Rbar}{\rho n}$ 精确对消：
\begin{align}
  \Delta_{\rho, n} &= -\rho\frac{\mathrm{d}\tau^2}{\mathrm{d}\rho}(\Rbar^2+v_R) - \frac{\rho\frac{\mathrm{d}\tau^2}{\mathrm{d}\rho}(\Rbar^2+v_R)}{kn} - \frac{(1+\tau^2)(\Rbar^2+v_R)}{kn} \nonumber \\
  &= \frac{\Rbar^2+v_R}{kn} \left[ -kn\rho\frac{\mathrm{d}\tau^2}{\mathrm{d}\rho} - \rho\frac{\mathrm{d}\tau^2}{\mathrm{d}\rho} - (1+\tau^2) \right] \nonumber \\
  &= \frac{\Rbar^2+v_R}{kn} \left[ -\rho\frac{\mathrm{d}\tau^2}{\mathrm{d}\rho}(kn + 1) - (1 + \tau^2) \right].
\end{align}
式~\eqref{eq:reversal-condition} 获证。

\paragraph{情况 1：制度型波动。}
当波动来自医院填报口径切换或法案调整时，系统的相对离散度 $\tau$ 取决于管理标准，不随基础报告率 $\rho$ 的提高而必然衰减，即 $\frac{\mathrm{d}\tau^2}{\mathrm{d}\rho} = 0$。代入式~\eqref{eq:reversal-condition}：
\begin{equation}
  \Delta_{\rho, n} = \frac{\Rbar^2+v_R}{kn} \big[ 0 - (1+\tau^2) \big] = -\frac{(1+\tau^2)(\Rbar^2+v_R)}{kn} < 0.
\end{equation}
差值恒严格为负。这意味着：扩大覆盖面 $n$ 带来的方差下降永远严格大于提高报告率 $\rho$。

\paragraph{情况 2：贝塔--二项抽样型波动。}
设每个周期的报告概率 $p_t$ 自身是由有效规模为 $\nu$ 的哨点样本抽样确定的，即满足共轭贝塔先验 $p_t \sim \Beta(\rho\nu, (1-\rho)\nu)$。根据贝塔分布的方差公式：
\begin{equation}
  \Var(p_t) = \frac{(\rho\nu)[(1-\rho)\nu]}{\nu^2(\nu + 1)} = \frac{\rho(1-\rho)}{\nu + 1}.
\end{equation}
计算相对方差 $\tau^2$ 及其对 $\rho$ 的微商：
\begin{align}
  \tau^2(\rho) &= \frac{\Var(p_t)}{\rho^2} = \frac{1-\rho}{\rho(\nu + 1)} = \frac{1}{\nu+1}\left( \frac{1}{\rho} - 1 \right), \\
  \frac{\mathrm{d}\tau^2}{\mathrm{d}\rho} &= -\frac{1}{\rho^2(\nu + 1)} < 0 \implies -\rho\frac{\mathrm{d}\tau^2}{\mathrm{d}\rho} = \frac{1}{\rho(\nu + 1)}.
\end{align}
同时计算常数项：
\begin{equation}
  1 + \tau^2 = 1 + \frac{1-\rho}{\rho(\nu + 1)} = \frac{\rho(\nu+1) + 1 - \rho}{\rho(\nu + 1)} = \frac{\rho\nu + 1}{\rho(\nu + 1)}.
\end{equation}
将微商项代入判别式括号内部：
\begin{align}
  -\rho\frac{\mathrm{d}\tau^2}{\mathrm{d}\rho}(kn + 1) - (1 + \tau^2) &= \frac{1}{\rho(\nu+1)}(kn + 1) - \frac{\rho\nu + 1}{\rho(\nu+1)} \nonumber \\
  &= \frac{(kn + 1) - (\rho\nu + 1)}{\rho(\nu + 1)} \nonumber \\
  &= \frac{kn - \rho\nu}{\rho(\nu + 1)}.
\end{align}
因此，提高报告率占优的充要条件是括号内为正：
\begin{equation}
  \Delta_{\rho, n} > 0 \iff kn - \rho\nu > 0 \iff kn > \rho\nu.
\end{equation}
两端同乘报告率 $\rho$ 并利用期望规模定义 $m = \rho n$：
\begin{equation}
  k(\rho n) > \rho^2 \nu \iff km > \rho^2 \nu \iff m > \frac{\rho^2 \nu}{k}.
\end{equation}
定理获证。
\end{proof}

\begin{derivationbox}{数值算例 7.1：决策反转阈值的实际量级}
\begin{itemize}
  \item \textbf{中等抽样哨点规模 ($\nu = 100$)：}
  设定流感基准参数 $\rho = 0.25, k = 0.35$：
  \begin{equation}
    m^* = \frac{0.25^2 \times 100}{0.35} = \frac{6.25}{0.35} \approx 17.86 \text{ 例}.
  \end{equation}
  当监测系统单期捕获病例数 $m > 17.9$ 例时，抽样波动的方差压制优势倒向提高报告率！
  \item \textbf{高精度哨点规模 ($\nu = 1{,}000$)：}
  \begin{equation}
    m^* = \frac{0.25^2 \times 1000}{0.35} \approx 178.6 \text{ 例}.
  \end{equation}
  系统必须维持在每周近 180 例以上的高负荷时，报告率的提升才会压倒人群覆盖的扩展。
\end{itemize}
\end{derivationbox}

\section{滞后分母加权与诊断性报告签名 $F(w, \psi)$}
\label{sec:part08}

\subsection{增长比统计量与分母回溯结构}

在时间序列被动监测中，流行病学家通常用当周报告病例数与前期报告病例数的加权和构造经验增长比统计量：
\begin{equation}
  G_t \equiv \frac{C_{t+1}}{D_t}, \qquad D_t \equiv \sum_{s=1}^K w_s C_{t+1-s}, \quad \sum_{s=1}^K w_s = 1.
  \label{eq:Gt-def}
\end{equation}
为了检验经验截距中的不可约误差究竟是由环境生物学波动（$v_R$）主导，还是由监测系统的报告率冲击（$v_H$）主导，我们考察改变滞后权重向量 $w$ 时截距的响应规律。

\subsection{卷积核函数 $F(w, \psi)$ 的严格推导}

设各周的真实发病在局部平稳，但监测系统存在平稳乘性报告冲击 $u_t \equiv p_t/\rho$。假设 $\{u_t\}$ 服从一阶自回归 AR(1) 过程：
\begin{equation}
  u_t - 1 = \psi (u_{t-1} - 1) + \xi_t, \qquad \xi_t \sim \text{iid}(0, \sigma_\xi^2), \quad |\psi| < 1.
\end{equation}
由此，报告冲击的理论自协方差函数为：
\begin{equation}
  \Cov(u_s, u_t) = \tau^2 \psi^{|s - t|}, \qquad \forall s, t.
\end{equation}

\begin{theorembox}{命题 4(vi) 的完全证明（报告签名的解析卷积核）}
若观测误差完全由平稳乘性报告冲击驱动（$v_R = 0$，且忽略微观二项抽样噪声），则增长统计量 $G_t$ 的相对预测误差方差在渐近一阶近似下正比于卷积核函数 $F(w, \psi)$：
\begin{equation}
  \Var\left( \frac{G_t}{\Rbar} \right) = \tau^2 F(w, \psi) + o(\tau^2),
  \label{eq:var-Gt-F}
\end{equation}
其中卷积核 $F(w, \psi)$ 严格由下式定义：
\begin{equation}
  \boxed{\, F(w, \psi) \equiv 1 - 2\sum_{s=1}^K w_s \psi^s + \sum_{s=1}^K \sum_{s'=1}^K w_s w_{s'} \psi^{|s - s'|} \,}
  \label{eq:Fw-kernel-proof}
\end{equation}
\end{theorembox}

\begin{proof}
在无计数离散噪声的连续极限下，报告病例数严格等于真实感染数乘以随机报告率：$C_t = \rho u_t I_t$。代入增长统计量：
\begin{equation}
  G_t = \frac{\rho u_{t+1} I_{t+1}}{\sum_s w_s \rho u_{t+1-s} I_{t+1-s}}.
\end{equation}
常数平均报告率 $\rho$ 在分子与分母中精确对消。在局部平稳增长或均值基准附近，$I_{t+1} / I_{t+1-s} \approx \Rbar^s$。为隔离报告冲击，考察无噪声基准 $G_t^0 \equiv I_{t+1} / \sum_s w_s I_{t+1-s} = \Rbar$。相对摄动比为：
\begin{equation}
  \frac{G_t}{G_t^0} = \frac{u_{t+1}}{\sum_{s=1}^K w_s u_{t+1-s}}.
\end{equation}
设微扰项 $\epsilon_t \equiv u_t - 1$，满足 $\E(\epsilon_t) = 0, \Var(\epsilon_t) = \tau^2, \Cov(\epsilon_s, \epsilon_t) = \tau^2 \psi^{|s-t|}$。
分子为 $1 + \epsilon_{t+1}$。分母加权和为：
\begin{equation}
  \sum_{s=1}^K w_s u_{t+1-s} = \sum_{s=1}^K w_s (1 + \epsilon_{t+1-s}) = \sum_{s=1}^K w_s + \sum_{s=1}^K w_s \epsilon_{t+1-s} = 1 + \bar{\epsilon}_w,
\end{equation}
其中记分母加权误差 $\bar{\epsilon}_w \equiv \sum_{s=1}^K w_s \epsilon_{t+1-s}$。对分母在 $\bar{\epsilon}_w = 0$ 处进行一阶一阶泰勒展开：
\begin{equation}
  \frac{1}{1 + \bar{\epsilon}_w} = 1 - \bar{\epsilon}_w + O(\bar{\epsilon}_w^2).
\end{equation}
因此，相对扰动比为：
\begin{align}
  \frac{G_t}{G_t^0} - 1 &= (1 + \epsilon_{t+1})(1 - \bar{\epsilon}_w + O(\bar{\epsilon}_w^2)) - 1 \nonumber \\
  &= 1 + \epsilon_{t+1} - \bar{\epsilon}_w + O(\tau^2) - 1 \nonumber \\
  &= \epsilon_{t+1} - \bar{\epsilon}_w + O(\tau^2) \nonumber \\
  &= \epsilon_{t+1} - \sum_{s=1}^K w_s \epsilon_{t+1-s} + O(\tau^2).
\end{align}
计算该线性组合的总体方差：
\begin{align}
  \Var\left( \epsilon_{t+1} - \sum_{s=1}^K w_s \epsilon_{t+1-s} \right) &= \Var(\epsilon_{t+1}) + \Var\left( \sum_{s=1}^K w_s \epsilon_{t+1-s} \right) \nonumber \\
  &\quad - 2\,\Cov\left( \epsilon_{t+1},\ \sum_{s=1}^K w_s \epsilon_{t+1-s} \right) \nonumber \\
  &= \Var(\epsilon_{t+1}) - 2 \sum_{s=1}^K w_s \Cov(\epsilon_{t+1}, \epsilon_{t+1-s}) \nonumber \\
  &\quad + \sum_{s=1}^K \sum_{s'=1}^K w_s w_{s'} \Cov(\epsilon_{t+1-s}, \epsilon_{t+1-s'}).
\end{align}
代入 AR(1) 自协方差公式 $\Cov(\epsilon_u, \epsilon_v) = \tau^2 \psi^{|u-v|}$：
\begin{itemize}
  \item 第一项：$\Var(\epsilon_{t+1}) = \tau^2$；
  \item 第二项：$\Cov(\epsilon_{t+1}, \epsilon_{t+1-s}) = \tau^2 \psi^{(t+1) - (t+1-s)} = \tau^2 \psi^s$；
  \item 第三项：$\Cov(\epsilon_{t+1-s}, \epsilon_{t+1-s'}) = \tau^2 \psi^{|(t+1-s) - (t+1-s')|} = \tau^2 \psi^{|s - s'|}$。
\end{itemize}
提取公因子 $\tau^2$：
\begin{equation}
  \Var\left( \frac{G_t}{G_t^0} - 1 \right) = \tau^2 \left[ 1 - 2\sum_{s=1}^K w_s \psi^s + \sum_{s=1}^K \sum_{s'=1}^K w_s w_{s'} \psi^{|s - s'|} \right].
\end{equation}
方括号内即为式~\eqref{eq:Fw-kernel-proof} 的 $F(w, \psi)$。证毕。
\end{proof}

\subsection{纯报告噪声下的截距比极端边界推导}

\begin{derivationbox}{推导 8.1：三种滞后权重下截距比可达范围的解析证明}
论文考虑了三种具体的滞后权重方案：
\begin{enumerate}
  \item 主分析规格：$w_{\text{primary}} = (0.65, 0.25, 0.10)$，平均滞后 $\bar{s} = 0.65\times 1 + 0.25\times 2 + 0.10\times 3 = 1.45$ 周；
  \item 单周规格：$w^{(1)} = (1, 0, 0)$，平均滞后 $\bar{s} = 1.0$ 周；
  \item 三周等权规格：$w^{(3)} = (1/3, 1/3, 1/3)$，平均滞后 $\bar{s} = 2.0$ 周。
\end{enumerate}
根据定理 8.1，相对主规格的截距比值为：
\begin{equation}
  q_1(\psi) \equiv \frac{F(w^{(1)}, \psi)}{F(w_{\text{primary}}, \psi)}, \qquad q_2(\psi) \equiv \frac{F(w^{(3)}, \psi)}{F(w_{\text{primary}}, \psi)}.
\end{equation}
下面求当自相关系数 $\psi \in [0, 1)$ 变化时，$q_1$ 与 $q_2$ 的解析上下确界：

\paragraph{极端 1：完全白噪声无记忆冲击（$\psi = 0$）。}
当 $\psi = 0$ 时，$\psi^s = 0$ 对所有 $s \ge 1$ 成立，且 $\psi^{|s-s'|} = \mathbb{I}(s = s')$。
此时卷积核完全退化为权重平方和：
\begin{equation}
  F(w, 0) = 1 - 0 + \sum_{s=1}^K w_s^2 = 1 + \sum_{s=1}^K w_s^2.
\end{equation}
代入三种权重的数值：
\begin{align}
  F(w_{\text{primary}}, 0) &= 1 + (0.65^2 + 0.25^2 + 0.10^2) = 1 + 0.4950 = 1.4950, \\
  F(w^{(1)}, 0) &= 1 + 1^2 = 2.0000, \\
  F(w^{(3)}, 0) &= 1 + 3 \times (1/3)^2 = 1 + 3 \times \frac{1}{9} = 1 + \frac{1}{3} \approx 1.3333.
\end{align}
由此精确计算出 $\psi = 0$ 时的截距比：
\begin{align}
  q_1(0) &= \frac{2.0000}{1.4950} \approx 1.3378 \approx 1.338, \\
  q_2(0) &= \frac{1.3333}{1.4950} \approx 0.8919 \approx 0.892.
\end{align}

\paragraph{极端 2：强自相关持续性冲击（$\psi \to 1^-$）。}
当 $\psi \to 1$ 时，由于 $\sum w_s = 1$，分子与分母均趋于 $1 - 2(1) + 1 = 0$。
应用洛必达法则（L'H\^opital's rule），展开 $\psi = 1 - \delta$ 当 $\delta \to 0$：
\begin{align}
  \psi^s &= (1 - \delta)^s = 1 - s\delta + \frac{s(s-1)}{2}\delta^2 + O(\delta^3), \\
  \psi^{|s-s'|} &= 1 - |s-s'|\delta + \frac{|s-s'|(|s-s'|-1)}{2}\delta^2 + O(\delta^3).
\end{align}
代入 $F(w, \psi)$：
\begin{itemize}
  \item 常数项：$1 - 2\sum w_s + \sum\sum w_s w_{s'} = 1 - 2(1) + (1)(1) = 0$；
  \item 一阶项 $\delta$：$-2\sum w_s(-s) - \sum\sum w_s w_{s'}|s-s'| = 2\bar{s} - \E|S - S'|$，其中 $S, S'$ 是独立服从权重分布 $w$ 的随机滞后。由于 $\E|S-S'| = 2\Cov(S, F(S)) > 0$，此项非零；
\end{itemize}
对各多项式精确展开求比值极限（亦经计算机代数系统精确符号求值）：
\begin{align}
  \lim_{\psi \to 1^-} q_1(\psi) &= 0.8833 \approx 0.883, \\
  \lim_{\psi \to 1^-} q_2(\psi) &= 1.3742 \approx 1.374.
\end{align}
由于 $q_1(\psi)$ 关于 $\psi$ 单调递减，$q_2(\psi)$ 关于 $\psi$ 单调递增，纯报告噪声所能解释的截距比区间被严格锁定在：
\begin{equation}
  \boxed{\, q_1 \in [0.883, 1.338], \qquad q_2 \in [0.892, 1.374] \,}
\end{equation}
\end{derivationbox}

\begin{insightbox}{实证辨识结论 8.1}
美国 CDC NHSN 流感住院实证估计得出的 95\% 置信区间为：
\begin{equation}
  q_1 = 0.688 \ [0.658, 0.727], \qquad q_2 = 1.539 \ [1.420, 1.656].
\end{equation}
实证区间呈现 $q_1 < 1 < q_2$（单周截距更低，长滞后截距更高），且两个区间与上述纯报告噪声的解析可达范围 $[0.883, 1.338]$ 和 $[0.892, 1.374]$ 不相交。
这表明实证结果与所考察的局部小扰动 AR(1) 纯报告冲击模型不相容。

需要说明的是：从式~\eqref{eq:Gt-def} 到线性微扰展开，严格推导应当使用有效滞后权重 $\tilde{w}_s \equiv \frac{w_s I_{t+1-s}}{\sum_j w_j I_{t+1-j}}$；此处近似使用固定名义权重 $w_s$，依赖于前序历史感染水平在回溯窗口内相对平稳的局部假设。
\end{insightbox}

\section{监测设计杠杆、时间自相关与预算优化}
\label{sec:part09}

\subsection{多期平均误差分解公式的完全推导}

\begin{theorembox}{命题 6 的完全证明（监测设计三大杠杆的误差分工）}
在假设 H1--H5 下，设环境序列自相关函数为 $\gamma_R(h) \equiv \Cov(R_t, R_{t+h})/v_R$，标准化报告冲击自相关为 $\gamma_u(h) \equiv \Cov(u_t, u_{t+h})/\tau^2$。采用 $T$ 期算术均值 $\bar{R}_T = \frac{1}{T}\sum_{t=1}^T \widehat{R}_t$ 估计跨期均值 $\Rbar$ 时，其均方误差严格服从：
\begin{equation}
  \E[(\bar{R}_T - \Rbar)^2] = \frac{1}{T^2}\sum_{s=1}^T \sum_{t=1}^T \Big\{ v_R \gamma_R(|s-t|) + \tau^2 \gamma_u(|s-t|)\big[\Rbar^2 + v_R \gamma_R(|s-t|)\big] \Big\} + \frac{A_H}{Tm}.
  \label{eq:alloc-decomp-general}
\end{equation}
\end{theorembox}

\begin{proof}
将估计误差拆解为抽样扰动与宏观潜在状态扰动：
\begin{equation}
  \bar{R}_T - \Rbar = \frac{1}{T}\sum_{t=1}^T \left( \frac{C_t}{m} - u_t R_t \right) + \frac{1}{T}\sum_{t=1}^T (u_t R_t - \Rbar) \equiv \Xi_1 + \Xi_2.
\end{equation}
\begin{enumerate}
  \item \textbf{正交性证明：}在给定整个潜在过程轨道 $\mathcal{G} \equiv (R_1, \dots, R_T, p_1, \dots, p_T)$ 的条件下，$\E(C_t / m \mid \mathcal{G}) = \frac{p_t n R_t}{\rho n} = u_t R_t$。因此 $\E(\Xi_1 \mid \mathcal{G}) = 0$。由重期望法则：
  \begin{equation}
    \E(\Xi_1 \Xi_2) = \E\big[ \E(\Xi_1 \Xi_2 \mid \mathcal{G}) \big] = \E\big[ \Xi_2 \cdot \E(\Xi_1 \mid \mathcal{G}) \big] = \E[\Xi_2 \cdot 0] = 0.
  \end{equation}
  故两项不相关，总均方误差为各自方差之和：$\E[(\bar{R}_T - \Rbar)^2] = \Var(\Xi_1) + \Var(\Xi_2)$。
  \item \textbf{抽样项 $\Var(\Xi_1)$：}由条件独立性假设 H6，各期的观测抽样在给定潜在状态下条件独立：
  \begin{equation}
    \Var(\Xi_1) = \frac{1}{T^2}\sum_{t=1}^T \E\big[\Var(C_t / m \mid R_t, p_t)\big].
  \end{equation}
  由命题 3 的推导，$\E[\Var(C_t / m \mid R_t, p_t)] = \frac{A_H}{m}$。因此：
  \begin{equation}
    \Var(\Xi_1) = \frac{1}{T^2} \sum_{t=1}^T \frac{A_H}{m} = \frac{T A_H}{T^2 m} = \frac{A_H}{Tm} = \frac{A_H}{B},
  \end{equation}
  其中 $B = Tm$ 是总观测病例量。
  \item \textbf{宏观项 $\Var(\Xi_2)$：}计算双重求和协方差：
  \begin{equation}
    \Var(\Xi_2) = \frac{1}{T^2}\sum_{s=1}^T \sum_{t=1}^T \Cov(u_s R_s, u_t R_t).
  \end{equation}
  利用假设 H5（报告率过程与传播过程相互独立）以及 $\E(u_t) = 1$：
  \begin{align}
    \Cov(u_s R_s, u_t R_t) &= \E(u_s u_t R_s R_t) - [\E(u_t R_t)]^2 \nonumber \\
    &= \E(u_s u_t) \, \E(R_s R_t) - \Rbar^2 \nonumber \\
    &= \big[ 1 + \tau^2 \gamma_u(|s-t|) \big] \big[ \Rbar^2 + v_R \gamma_R(|s-t|) \big] - \Rbar^2 \nonumber \\
    &= \Rbar^2 + v_R \gamma_R(|s-t|) + \tau^2 \gamma_u(|s-t|)\big[\Rbar^2 + v_R \gamma_R(|s-t|)\big] - \Rbar^2 \nonumber \\
    &= v_R \gamma_R(|s-t|) + \tau^2 \gamma_u(|s-t|)\big[\Rbar^2 + v_R \gamma_R(|s-t|)\big].
  \end{align}
\end{enumerate}
将两部分合并，即完全得出式~\eqref{eq:alloc-decomp-general}。
\end{proof}

\subsection{独立环境下的成本优化与 Neyman 费用分配}

当环境与报告在各期相互独立（$\gamma_R(h) = \gamma_u(h) = \mathbb{I}(h=0)$），且报告标准化（$\tau = 0$）时，多期平均方差简化为式~\eqref{eq:alloc-indep}：
\begin{equation}
  V_T(m) = \frac{v_R}{T} + \frac{A}{Tm}.
  \label{eq:alloc-indep}
\end{equation}

\begin{theorembox}{推导 9.1：预算约束下拉格朗日极值证明}
设每增加一个观测时期的固定场地/网络运维成本为 $c_0$，每追踪采集一例病例的边际成本为 $c_1$。总经费预算为 $\mathcal{B}$，预算约束为：
\begin{equation}
  T(c_0 + c_1 m) = \mathcal{B} \implies T = \frac{\mathcal{B}}{c_0 + c_1 m}.
\end{equation}
将 $T$ 代入方差目标函数：
\begin{equation}
  f(m) = \frac{v_R + A/m}{\frac{\mathcal{B}}{c_0 + c_1 m}} = \frac{1}{\mathcal{B}} (c_0 + c_1 m)\left( v_R + \frac{A}{m} \right) = \frac{1}{\mathcal{B}} \left( c_0 v_R + c_1 A + c_1 v_R m + \frac{c_0 A}{m} \right).
\end{equation}
对单期规模 $m$ 求一阶导数：
\begin{equation}
  f'(m) = \frac{1}{\mathcal{B}} \left( c_1 v_R - \frac{c_0 A}{m^2} \right).
\end{equation}
令 $f'(m) = 0$：
\begin{equation}
  c_1 v_R = \frac{c_0 A}{m^2} \iff m^2 = \frac{c_0 A}{c_1 v_R} = \frac{c_0}{c_1} \cdot \frac{A}{v_R} = \frac{c_0}{c_1} \mx.
\end{equation}
由此得出唯一驻点：
\begin{equation}
  \boxed{\, m^* = \sqrt{\frac{c_0}{c_1} \mx} \,}
  \label{eq:opt-m-neyman}
\end{equation}
求二阶导数核验凸性：
\begin{equation}
  f''(m) = \frac{1}{\mathcal{B}} \left( \frac{2 c_0 A}{m^3} \right) > 0, \quad \forall m > 0.
\end{equation}
目标函数在正实轴上严格凸，驻点必为全局唯一最小值。将 $m^*$ 代回目标函数：
\begin{align}
  f(m^*) &= \frac{1}{\mathcal{B}} \left( c_0 v_R + c_1 A + 2\sqrt{c_1 v_R m^* \cdot \frac{c_0 A}{m^*}} \right) \nonumber \\
  &= \frac{1}{\mathcal{B}} \left( c_0 v_R + c_1 A + 2\sqrt{c_0 c_1 v_R A} \right) \nonumber \\
  &= \frac{(\sqrt{c_0 v_R} + \sqrt{c_1 A})^2}{\mathcal{B}}.
\end{align}
\end{theorembox}

\begin{insightbox}{理论联系 9.1：两阶段抽样费用最优分配的同构性}
式~\eqref{eq:opt-m-neyman} 在数学结构上与经典抽样调查理论中的 Neyman--Cochran 两阶段抽样费用最优分配定理 \citep{cochran1977sampling} \textbf{完全同构}。
最优单期样本量取决于固定成本与边际成本之比 $\sqrt{c_0/c_1}$ 以及方差平衡规模 $\sqrt{\mx}$ 的几何均值。
\end{insightbox}

\subsection{时间自相关 AR(1) 环境与关键预算阈值 $\mathcal{B}^\dagger$}

当真实传播条件存在时间记忆性（如天气和季节在相邻周保持连贯），设 $R_t$ 服从平稳高斯 AR(1) 过程：$\Cov(R_s, R_t) = v_R \phi^{|s-t|}$（$\phi \in (0, 1)$）。

\begin{derivationbox}{推导 9.2：自回归环境下的总方差与有效样本量}
计算均值方差项 $\Var(\bar{R}_T) = \frac{v_R}{T^2}\sum_{s=1}^T \sum_{t=1}^T \phi^{|s-t|}$：
将方阵沿主对角线与次对角线求和：
\begin{equation}
  \sum_{s=1}^T \sum_{t=1}^T \phi^{|s-t|} = T + 2\sum_{h=1}^{T-1} (T - h) \phi^h = T \left[ 1 + 2\sum_{h=1}^{T-1} \left(1 - \frac{h}{T}\right)\phi^h \right] \equiv T \kappa_T(\phi).
\end{equation}
利用幂级数求和公式计算内层求和：
\begin{align}
  \sum_{h=1}^{T-1} \phi^h &= \frac{\phi(1 - \phi^{T-1})}{1 - \phi}, \\
  \sum_{h=1}^{T-1} h \phi^h &= \phi \frac{\mathrm{d}}{\mathrm{d}\phi} \sum_{h=1}^{T-1} \phi^h = \frac{\phi - T\phi^T + (T-1)\phi^{T+1}}{(1 - \phi)^2}.
\end{align}
通分相消整理，得出精确闭式：
\begin{equation}
  T^2 \Var(\bar{R}_T) = v_R \left[ \frac{T(1+\phi)}{1-\phi} - \frac{2\phi(1-\phi^T)}{(1-\phi)^2} \right].
\end{equation}
定义单位周期有效相对环境方差因子 $e(T) \equiv \frac{\kappa_T(\phi)}{T}$。当预算约束 $T(c_0 + c_1 m) = \mathcal{B}$ 给出 $m = \frac{\mathcal{B} - c_0 T}{c_1 T}$ 时，目标函数关于连续变量期数 $T$ 表达为：
\begin{equation}
  f_\phi(T) = v_R e(T) + \frac{A c_1}{\mathcal{B} - c_0 T}, \qquad 1 \le T < \frac{\mathcal{B}}{c_0}.
\end{equation}
定义边际检验核函数 $g_T(\phi) \equiv -T^2 e'(T)$：
\begin{equation}
  g_T(\phi) = \frac{1+\phi}{1-\phi} - \frac{2\phi}{(1-\phi)^2} \left[ \frac{2(1-\phi^T)}{T} + \phi^T \ln\phi \right].
\end{equation}
经求导可证：$g_T(\phi)$ 关于 $T$ 处处严格单调递增，且在端点满足：
\begin{equation}
  g_1(\phi) < 1, \qquad \lim_{T \to \infty} g_T(\phi) = \frac{1+\phi}{1-\phi} > 1.
\end{equation}
由介值定理，$g_T(\phi) = 1$ 存在唯一的根 $T^\dagger(\phi)$。
对应地，存在\textbf{临界预算阈值}：
\begin{equation}
  \mathcal{B}^\dagger \equiv T^\dagger (c_0 + c_1 m_0^*), \qquad \text{其中 } m_0^* = \sqrt{\frac{c_0}{c_1}\mx}.
\end{equation}
\begin{itemize}
  \item \textbf{小预算体系 ($\mathcal{B} < \mathcal{B}^\dagger$)：}解满足 $T^* < T^\dagger \implies g_{T^*} < 1 \implies m^* > m_0^*$。
  \textbf{最优策略是减少期数，集中预算做大单期规模！}
  \item \textbf{大预算体系 ($\mathcal{B} > \mathcal{B}^\dagger$)：}解满足 $T^* > T^\dagger \implies g_{T^*} > 1 \implies m^* < m_0^*$。
  \textbf{最优策略是缩减单期规模，将资金用于增加独立观测期数！}
\end{itemize}
\end{derivationbox}

\subsection{空间多辖区协同汇总（Spatial Pooling）}

\begin{theorembox}{注 2 的完全证明（空间池化下的最优规模）}
设协同监测网络涵盖 $J$ 个平行辖区，辖区两两之间环境波动的平均相关系数为 $\bar{c} \in [0, 1]$。总预算为 $\mathcal{B} = J(c_0 + c_1 m)$。则最优单辖区规模为：
\begin{equation}
  m^* = \sqrt{\frac{c_0 A}{c_1 v_R (1 - \bar{c})}}.
  \label{eq:opt-spatial}
\end{equation}
\end{theorembox}

\begin{proof}
根据式~\eqref{eq:pooling}，等权汇总 $J$ 个辖区后的总均方误差为：
\begin{equation}
  \text{MSE}_J(m) = v_R \bar{c} + \frac{v_R(1 - \bar{c})}{J} + \frac{A}{Jm}.
  \label{eq:pooling}
\end{equation}
代入预算约束 $J = \frac{\mathcal{B}}{c_0 + c_1 m}$：
\begin{equation}
  \text{MSE}(m) = v_R \bar{c} + \frac{c_0 + c_1 m}{\mathcal{B}} \left[ v_R(1 - \bar{c}) + \frac{A}{m} \right].
\end{equation}
展开方括号项：
\begin{equation}
  \text{MSE}(m) = v_R \bar{c} + \frac{1}{\mathcal{B}} \left[ c_0 v_R(1-\bar{c}) + c_1 A + c_1 v_R(1-\bar{c})m + \frac{c_0 A}{m} \right].
\end{equation}
对单期单辖区规模 $m$ 求导令其等于零：
\begin{equation}
  c_1 v_R (1 - \bar{c}) - \frac{c_0 A}{m^2} = 0 \implies m^2 = \frac{c_0 A}{c_1 v_R (1 - \bar{c})}.
\end{equation}
开平方即得式~\eqref{eq:opt-spatial}。证毕。
\end{proof}

\section{环境方差估计量 $\hat{v}_R$ 的有限样本精度与信息界限}
\label{sec:part10}

\subsection{修正矩估计量 $\hat{v}_R$ 的严格代数构造}

由于环境方差 $v_R$ 自身必须从有限期数和有限规模的报告数据中估计，我们推导其修正矩估计量。

\begin{derivationbox}{推导 10.1：解二阶总体矩方程得出估计量}
设单期标准化观测序列为 $Y_t \equiv C_t / m$。各期规模同为 $m = \rho n$。
由式~\eqref{eq:var-closed}，未标准化的总报告数边际方差为：
\begin{align}
  \Var(C_t) &= m^2 V(m) = m^2 \left( v_R + \frac{A}{m} \right) = m^2 v_R + m A \nonumber \\
  &= m^2 v_R + m \left[ \Rbar + \frac{\rho(\Rbar^2 + v_R)}{k} \right] \nonumber \\
  &= m^2 v_R + m \Rbar + \frac{\rho m \Rbar^2}{k} + \frac{\rho m}{k} v_R.
\end{align}
合并含待估参数 $v_R$ 的项：
\begin{equation}
  \Var(C_t) = v_R \left( m^2 + \frac{\rho m}{k} \right) + m \Rbar + \frac{\rho m \Rbar^2}{k} = v_R \left( \frac{k m^2 + \rho m}{k} \right) + m \Rbar + \frac{\rho m \Rbar^2}{k}.
\end{equation}
用样本二阶中心矩 $S_C^2 \equiv \frac{1}{T-1}\sum_{t=1}^T (C_t - \bar{C})^2$ 替换 $\Var(C_t)$，用样本均值 $\bar{C}$ 替换 $m\Rbar$（即 $\Rbar = \bar{C}/m$）：
\begin{align}
  S_C^2 &= v_R \left( \frac{k m^2 + \rho m}{k} \right) + \bar{C} + \frac{\rho m}{k} \left( \frac{\bar{C}}{m} \right)^2 \nonumber \\
  &= v_R \left( \frac{k m^2 + \rho m}{k} \right) + \bar{C} + \frac{\rho \bar{C}^2}{km}.
\end{align}
移项解 $v_R$：
\begin{equation}
  v_R \left( \frac{k m^2 + \rho m}{k} \right) = S_C^2 - \bar{C} - \frac{\rho \bar{C}^2}{km} = \frac{k m S_C^2 - k m \bar{C} - \rho \bar{C}^2}{km}.
\end{equation}
两侧同乘以 $k$，同除以 $m^2(km + \rho)$：
\begin{equation}
  \hat{v}_R = \frac{k m S_C^2 - k m \bar{C} - \rho \bar{C}^2}{m^2 (km + \rho)}.
  \label{eq:vhat-raw}
\end{equation}
若用标准化样本统计量 $\bar{Y} = \bar{C}/m, S_Y^2 = S_C^2 / m^2$ 表示：
\begin{equation}
  \hat{v}_R = \frac{S_Y^2 - \left( \frac{\bar{Y}}{m} + \frac{\rho \bar{Y}^2}{km} \right)}{1 + \frac{\rho}{km}} = \frac{S_Y^2 - \varpi(\bar{Y})}{1 + \frac{\rho}{km}}, \qquad \varpi(y) \equiv \frac{y}{m} + \frac{\rho y^2}{km}.
  \label{eq:r3-vhat}
\end{equation}
式~\eqref{eq:r3-vhat} 获证。
\end{derivationbox}

\subsection{有限样本方差公式与极限平台 $\Omega_\infty / T$}

\begin{theorembox}{补充命题 S1 的证明（环境方差估计精度的极限平台）}
在独立同分布环境下，当单期规模 $m \to \infty$ 时，估计方差 $T \Var(\hat{v}_R)$ 不会趋向于 0，而是收敛至不可跨越的渐近平台：
\begin{equation}
  \lim_{m \to \infty} T \Var(\hat{v}_R) = \Omega_\infty \equiv \E(R_t - \Rbar)^4 - v_R^2 = \Var\big((R_t - \Rbar)^2\big).
  \label{eq:omega-inf-def}
\end{equation}
若 $R_t \sim \text{Gamma}(\xi, \vartheta)$，则极限平台具有严格代数闭式：
\begin{equation}
  \Omega_\infty = v_R^2 \left( 2 + \frac{6 v_R}{\Rbar^2} \right).
  \label{eq:omega-gamma-closed}
\end{equation}
\end{theorembox}

\begin{proof}
根据抽样方差的经典有限样本四阶矩理论 \citep{casella2002statistical}，样本方差的精确方差为：
\begin{equation}
  \Var(S_Y^2) = \frac{\mu_{4Y}}{T} - \frac{T - 3}{T(T - 1)} \sigma_Y^4,
\end{equation}
其中 $\mu_{4Y} \equiv \E(Y_t - \E Y_t)^4$ 为四阶中心矩，$\sigma_Y^2 = \Var(Y_t)$。
将 $\hat{v}_R = \frac{S_Y^2 - \varpi(\bar{Y})}{1 + \rho/(km)}$ 在 $\bar{Y} = \Rbar$ 处进行一阶一阶泰勒展开：
\begin{equation}
  \varpi(\bar{Y}) = \varpi(\Rbar) + \varpi'(\Rbar)(\bar{Y} - \Rbar) + O\left(\frac{1}{mT}\right), \qquad \varpi'(y) = \frac{1 + 2\rho y/k}{m}.
\end{equation}
代入计算协方差与方差组合：
\begin{equation}
  \Var(\hat{v}_R) = \frac{1}{\left(1 + \frac{\rho}{km}\right)^2} \left[ \frac{\mu_{4Y}}{T} - \frac{T - 3}{T(T - 1)} \sigma_Y^4 - 2\varpi'\frac{\mu_{3Y}}{T} + (\varpi')^2 \frac{\sigma_Y^2}{T} \right] \left(1 + O\left(\frac{1}{mT}\right)\right).
  \label{eq:r3-omega}
\end{equation}
利用阶乘矩已知 $\Cov(S_Y^2, \bar{Y}) = \frac{\mu_{3Y}}{T}, \Var(\bar{Y}) = \frac{\sigma_Y^2}{T}$，得出式~\eqref{eq:r3-omega}。

当 $m \to \infty$ 时，$Y_t \equiv C_t/m \xrightarrow{a.s.} R_t$，$\varpi' \to 0$，$1 + \frac{\rho}{km} \to 1$。
四阶矩收敛至真实环境的四阶矩 $\mu_{4Y} \to \mu_{4R}$，中心方差收敛至 $\sigma_Y^2 \to v_R$。因此：
\begin{equation}
  \lim_{m \to \infty} T \Var(\hat{v}_R) = \mu_{4R} - \frac{T-3}{T-1} v_R^2 \xrightarrow{T \to \infty} \mu_{4R} - v_R^2 \equiv \Omega_\infty.
\end{equation}
若 $R_t \sim \text{Gamma}(\xi, \vartheta)$，其均值为 $\xi\vartheta = \Rbar$，方差为 $\xi\vartheta^2 = v_R$。
解出形状参数与尺度参数：
\begin{equation}
  \xi = \frac{\Rbar^2}{v_R}, \qquad \vartheta = \frac{v_R}{\Rbar}.
\end{equation}
伽马分布的四阶中心矩为：
\begin{equation}
  \mu_{4R} = 3\xi(\xi + 2)\vartheta^4 = 3\left( \frac{\Rbar^2}{v_R} \right)\left( \frac{\Rbar^2}{v_R} + 2 \right)\left( \frac{v_R}{\Rbar} \right)^4 = 3 v_R^2 \left( 1 + \frac{2 v_R}{\Rbar^2} \right).
\end{equation}
代入 $\Omega_\infty$ 的定义：
\begin{equation}
  \Omega_\infty = \mu_{4R} - v_R^2 = 3 v_R^2 \left( 1 + \frac{2 v_R}{\Rbar^2} \right) - v_R^2 = 2 v_R^2 + \frac{6 v_R^3}{\Rbar^2} = v_R^2 \left( 2 + \frac{6 v_R}{\Rbar^2} \right).
\end{equation}
式~\eqref{eq:omega-gamma-closed} 获证。
\end{proof}

\subsection{多流行季联合推断与季节方差下界（Corollary 2 证明）}

在多年季节性传染病（如流感、RSV）中，环境方差包含跨年份的宏观季节效应。设每个流行季 $s \in \{1, \dots, S\}$ 具有不可约的季节特异性冲击：
\begin{equation}
  R_{s, t} = \Rbar + \eta_s + \omega_{s, t}, \qquad \eta_s \overset{\text{iid}}{\sim} \mathcal{N}(0, v_S).
\end{equation}

\begin{theorembox}{推论 2 的完全证明（季节方差估计精度的下界界限）}
无论每个流行季内部覆盖多少个辖区（$J \to \infty$）、采集多少周的序列（$T \to \infty$）、捕获多少例病例（$m \to \infty$），季层面方差 $v_S$ 的任意无偏估计量的相对标准误差（RSE）均存在不可逾越的统计下界：
\begin{equation}
  \text{RSE}(\hat{v}_S) \equiv \frac{\sqrt{\Var(\hat{v}_S)}}{v_S} \ge \sqrt{\frac{2}{S}}.
  \label{eq:rse-bound}
\end{equation}
当季节全局均值未知，采用样本方差估计时，下界提升为 $\sqrt{\frac{2}{S - 1}}$。
\end{theorembox}

\begin{proof}
构建系统参数、潜变量与全部微观观测数据的马尔可夫链：
\begin{equation}
  v_S \longrightarrow (\eta_1, \dots, \eta_S) \longrightarrow \{R_{s,t}\} \longrightarrow \{C_{s,t}\}.
\end{equation}
根据 Fisher 信息量的数据处理不等式（Data Processing Inequality）：
\begin{equation}
  \mathcal{I}_{\text{data}}(v_S) \le \mathcal{I}_{\eta_{1:S}}(v_S).
\end{equation}
右侧是在完美、无噪声直接观测到 $S$ 个季节真实潜在冲击 $\eta_1, \dots, \eta_S \overset{\text{iid}}{\sim} \mathcal{N}(0, v_S)$ 时的理想信息量。
由于零均值正态随机样本关于方差参数的对数似然为：
\begin{equation}
  \ell(v_S) = -\frac{S}{2}\ln(2\pi) - \frac{S}{2}\ln(v_S) - \frac{1}{2v_S}\sum_{s=1}^S \eta_s^2,
\end{equation}
求二阶导数并取负期望：
\begin{equation}
  \frac{\partial^2 \ell}{\partial v_S^2} = \frac{S}{2 v_S^2} - \frac{1}{v_S^3}\sum_{s=1}^S \eta_s^2 \implies \mathcal{I}_{\eta_{1:S}}(v_S) = -\E\left( \frac{\partial^2 \ell}{\partial v_S^2} \right) = \frac{S}{2 v_S^2}.
\end{equation}
因此，全样本关于 $v_S$ 的 Fisher 信息量绝对不可能超过 $\frac{S}{2 v_S^2}$。
由 Cramér--Rao 下界，对任意无偏估计量 $\hat{v}_S$：
\begin{equation}
  \Var(\hat{v}_S) \ge \frac{1}{\mathcal{I}_{\text{data}}(v_S)} \ge \frac{1}{\frac{S}{2 v_S^2}} = \frac{2 v_S^2}{S}.
\end{equation}
两端开平方并除以 $v_S$：
\begin{equation}
  \frac{\sqrt{\Var(\hat{v}_S)}}{v_S} \ge \sqrt{\frac{2}{S}}.
\end{equation}
若季节宏观均值未知，样本方差 $S_\eta^2 = \frac{1}{S-1}\sum_{s=1}^S (\eta_s - \bar{\eta})^2$ 的自由度为 $S-1$，方差等于 $\frac{2 v_S^2}{S-1}$，下界相应为 $\sqrt{\frac{2}{S-1}}$。证毕。
\end{proof}

\begin{derivationbox}{数值评价 10.1：多季流行病监测的精度瓶颈}
代入真实世界可获得的流行季数量：
\begin{itemize}
  \item $S = 3$ 季（本文主分析样本）：$\sqrt{2/S} = \sqrt{2/3} \approx 0.816 \approx 0.82$；均值未知时 $\sqrt{2/(3-1)} = 1.00$。这意味着相对标准差高达 100\%，置信区间极宽，必然跨越 0。
  \item $S = 4$ 季（纳入 2025--2026 样本外季）：$\sqrt{2/4} \approx 0.707 \approx 0.71$；均值未知时 $\sqrt{2/3} \approx 0.82$。
\end{itemize}
这在数学上彻底解释了：为什么无论纳入多少个州、多么庞大的住院病例，季随机效应估计的方差区间依然极为宽阔。这种不确定性是由大自然的季节轮回次数决定的，无法依靠在单季内堆砌病例来解决。
\end{derivationbox}

\section{假设检验、单调似然比与预警功效曲线}
\label{sec:part11}

\subsection{暴发预警的统计检验构造与单调似然比}

在传染病监测预警中（估计目标 1），决策者最核心的关切不是参数的渐近均方误差，而是及时判断疫情是否处于扩张阶段：
\begin{equation}
  H_0\colon R_t \le 1 \quad \text{versus} \quad H_1\colon R_t > 1.
  \label{eq:hypo-def}
\end{equation}

\begin{theorembox}{定理 11.1（负二项分布族的严格单调似然比性质）}
在假设 H1--H4 下，给定祖代规模 $n$，观测病例数 $C_t$ 的条件分布族关于参数 $R_t$ 具有严格单调似然比（Monotone Likelihood Ratio, MLR）。
\end{theorembox}

\begin{proof}
任取两个参数值 $R_2 > R_1 > 0$。根据式~\eqref{eq:der-ct-dist}，观测值 $C_t = c$ 在两个参数下的概率质量函数之比为：
\begin{align}
  \Lambda(c; R_2, R_1) &\equiv \frac{P(C_t = c \mid R_2)}{P(C_t = c \mid R_1)} \nonumber \\
  &= \frac{\frac{\Gamma(c+nk)}{c!\,\Gamma(nk)} \left(\frac{k}{k+\rho R_2}\right)^{nk} \left(\frac{\rho R_2}{k+\rho R_2}\right)^c}{\frac{\Gamma(c+nk)}{c!\,\Gamma(nk)} \left(\frac{k}{k+\rho R_1}\right)^{nk} \left(\frac{\rho R_1}{k+\rho R_1}\right)^c} \nonumber \\
  &= \left( \frac{k+\rho R_1}{k+\rho R_2} \right)^{nk} \left[ \frac{\rho R_2 / (k+\rho R_2)}{\rho R_1 / (k+\rho R_1)} \right]^c \nonumber \\
  &= \left( \frac{k+\rho R_1}{k+\rho R_2} \right)^{nk} \left[ \frac{R_2(k+\rho R_1)}{R_1(k+\rho R_2)} \right]^c.
  \label{eq:lr-ratio}
\end{align}
考察以 $c$ 为指数的底数项：
\begin{equation}
  \Theta(R_2, R_1) \equiv \frac{R_2(k+\rho R_1)}{R_1(k+\rho R_2)} = \frac{k R_2 + \rho R_1 R_2}{k R_1 + \rho R_1 R_2}.
\end{equation}
由于 $R_2 > R_1$ 且 $k > 0$，比较分子与分母：
\begin{equation}
  (k R_2 + \rho R_1 R_2) - (k R_1 + \rho R_1 R_2) = k(R_2 - R_1) > 0.
\end{equation}
因此分子严格大于分母，底数满足：
\begin{equation}
  \Theta(R_2, R_1) > 1.
\end{equation}
由于底数大于 1，比值函数 $\Lambda(c; R_2, R_1)$ 是统计量 $c$ 的\textbf{严格单调递增函数}。
证毕。
\end{proof}

\subsection{Karlin--Rubin 定理与一致最优（UMP）精确检验}

根据数理统计中的 Karlin--Rubin 定理 \citep{karlin1956theory, lehmann2005testing}，对于具有单调似然比性质的单参数分布族，单侧假设检验的一致最优势检验（Uniformly Most Powerful, UMP）必为临界域形式。

对于离散型负二项分布，由于概率累积存在阶梯跳跃，标准 Karlin--Rubin UMP 检验在理论上包含边界随机化：当 $C_t > c$ 时拒绝，当 $C_t = c$ 时以概率 $\gamma \in [0, 1)$ 随机化拒绝，从而使得实际水平严格等于名义水平 $\alpha$。在公共卫生监测工程实践中，为了避免人为抛硬币随机化导致的判断不确定性，通常采用确定性的保守非随机化检验规则。

\begin{derivationbox}{推导 11.1：临界值与精确功效函数}
设名义显著性水平为 $\alpha \in (0, 1)$。定义临界值 $c_\alpha$ 为使得零假设临界点 $R_t = 1$ 下第一类错误概率不超过 $\alpha$ 的最小整数：
\begin{equation}
  c_\alpha \equiv \min \left\{ c \in \mathbb{N} \ \Big|\ P_{R_t = 1}(C_t \ge c) \le \alpha \right\}.
\end{equation}
本文采用非随机化确定性规则：当且仅当 $C_t \ge c_\alpha$ 时拒绝原假设 $H_0$。
\begin{enumerate}
  \item \textbf{实际水平（Size）：}由于功效函数在单调似然比族下关于 $R_t$ 单调非减，原假设区域内的最大犯错概率在边界 $R_t = 1$ 处取得：
  \begin{equation}
    \alpha^* \equiv \sup_{R_t \le 1} P_{R_t}(C_t \ge c_\alpha) = P_{R_t = 1}(C_t \ge c_\alpha) \le \alpha.
  \end{equation}
  非随机化检验 $\{C_t \ge c_\alpha\}$ 对应随机化检验在 $c = c_\alpha - 1, \gamma = 1$ 的特例，因此它在其实际达到的显著性水平 $\alpha^*$ 下是一致最优（UMP）的；若在边界 $c_\alpha - 1$ 处允许随机化，则可使实际水平严格等于名义水平 $\alpha$ 并取得更高功效。
  \item \textbf{精确功效（Power）：}当真实传播发生暴发，即备择假设点 $R_t = 1 + \delta$（$\delta > 0$）时，观测病例数服从负二项分布：
  \begin{equation}
    C_t \sim \NB\!\left(nk,\ \frac{k}{k + \rho(1+\delta)}\right).
  \end{equation}
  其精确检测功效定义为拒绝原假设的概率，由负二项分布累积分布函数 $F_{\NB}$ 直接给出：
  \begin{equation}
    \beta(n, \delta) \equiv P_{R_t = 1+\delta}(C_t \ge c_\alpha) = 1 - F_{\NB}\left( c_\alpha - 1;\ nk,\ \frac{k}{k + \rho(1+\delta)} \right).
    \label{eq:exact-power-proof}
  \end{equation}
\end{enumerate}
\end{derivationbox}

\subsection{检验功效锯齿震荡与 Wald 检验的尺寸膨胀}

在实际工程应用中，许多流行病学软件常使用大样本正态近似的 Wald 检验统计量：
\begin{equation}
  Z \equiv \frac{\widehat{R}_t - 1}{\sqrt{V_0}} \sim \mathcal{N}(0, 1), \qquad V_0 \equiv \Var(\widehat{R}_t \mid R_t = 1) = \frac{1}{\rho n} + \frac{1}{kn}.
\end{equation}

\begin{insightbox}{数值病态洞见 11.1：Wald 检验严重虚假报警机制}
由于负二项分布在样本量较小及存在超传播（小 $k$）时呈现极端的右偏重尾特征，正态近似在尾部严重失真：
\begin{itemize}
  \item 在基准参数 $k = 0.1, m = 5$ 例的小样本下，名义水平为 5\% 的 Wald 检验，其实际第一类错误概率高达 \textbf{7.81\%}；
  \item 在 $m < 40$ 的区间内，不同超传播强度 $k = 0.1, 0.35, 1.0$ 下 Wald 检验的峰值假阳性率分别达到 \textbf{8.03\%}、\textbf{8.75\%} 和 \textbf{10.40\%}，超出安全标准一倍以上；
  \item 相比之下，基于式~\eqref{eq:exact-power-proof} 的精确检验在全参数网格上的最大实际检验水平为 \textbf{0.049997}，严格被 5\% 上界控制。
\end{itemize}

此外，由于临界值 $c_\alpha$ 随样本规模 $n$ 呈现阶梯状整数跳跃，精确功效曲线 $\beta(n, \delta)$ 关于 $n$ 表现出剧烈的\textbf{锯齿状震荡（Serration）}：
例如，在 $k = 0.35, \delta = 0.3$ 时，当祖代规模 $n$ 从 400 增加至 699 的过程中，由于每当临界点 $c_\alpha$ 增加 1，拒绝阈值突然上抬，导致检测功效在此过程中\textbf{逆向下降了 82 次}！

因此，监测系统必须区分两个规模标准：
\begin{enumerate}
  \item 首次达到 80\% 功效的脆弱规模；
  \item 越过最后一个锯齿跌落、此后稳定保持在 80\% 功效之上的鲁棒规模。
\end{enumerate}
在 $\rho = 0.25, k = 0.35$ 下，检出 15\% 暴发增长（$\delta = 0.15$）所需的两种祖代规模分别为 2{,}071 人与 2{,}107 人，对应预期病例数 $m \approx 527$ 例，约为方差平衡规模 $\mx = 53.1$ 的整整十倍。
\end{insightbox}

\section{动态更新过程与分母噪声二阶 Delta 展开}
\label{sec:part12}

\subsection{更新方程与感染压力不确定性分解}

在被动时间序列监测中，真实的祖代感染人数 $n$ 是不可观测的潜变量。根据流行病学更新方程（Renewal equation）\citep{fraser2007estimating, cori2013new}，第 $t+1$ 周产生的新发感染由历史上所有前序世代的发病共同激活，其宏观真实感染压力定义为：
\begin{equation}
  \Lambda_{t+1} \equiv \sum_{s=1}^\infty w_s I_{t+1-s},
\end{equation}
其中 $w_s$ 为离散化的世代间隔概率分布（满足 $\sum_s w_s = 1$）。由于真实发病数 $I_{t+1-s}$ 存在二项漏报，历史信息集 $\mathcal{F}_t$ 所能掌握的仅有报告病例的加权和 $D_t = \sum_s w_s C_{t+1-s}$。

\begin{theorembox}{定理 12.1（时间序列下感染压力不确定性的方差膨胀定理）}
设 $m \equiv \rho\,\E(\Lambda_{t+1} \mid \mathcal{F}_t)$ 为预期报告感染压力。条件在历史信息集 $\mathcal{F}_t$ 下，报告病例数 $C_{t+1}$ 的相对预测方差严格服从三项分解：
\begin{equation}
  \frac{\Var(C_{t+1} \mid \mathcal{F}_t)}{\big[\rho\,\E(\Lambda_{t+1} \mid \mathcal{F}_t)\big]^2} = v_R + \frac{A}{m} + \CV^2(\Lambda_{t+1} \mid \mathcal{F}_t)\,(\Rbar^2 + v_R),
  \label{eq:info-decomp-der}
\end{equation}
其中 $\CV(\Lambda_{t+1} \mid \mathcal{F}_t) \equiv \frac{\sqrt{\Var(\Lambda_{t+1} \mid \mathcal{F}_t)}}{\E(\Lambda_{t+1} \mid \mathcal{F}_t)}$ 为给定报告数据下潜在真实感染压力的变异系数。
\end{theorembox}

\begin{proof}
记 $\bar{\Lambda} \equiv \E(\Lambda_{t+1} \mid \mathcal{F}_t)$，则 $m = \rho \bar{\Lambda}$。所有期望与方差均在条件信息集 $\mathcal{F}_t$ 下计算。
第一层条件化：给定局部感染压力 $\Lambda_{t+1}$ 与环境 $R_{t+1}$。由定理 2.1，观测数服从负二项分布：
\begin{equation}
  C_{t+1} \mid \Lambda_{t+1}, R_{t+1} \sim \NB\!\left( k \Lambda_{t+1},\ \frac{k}{k + \rho R_{t+1}} \right).
\end{equation}
其条件期望与方差为：
\begin{align}
  \E(C_{t+1} \mid \Lambda_{t+1}, R_{t+1}) &= \rho \Lambda_{t+1} R_{t+1}, \\
  \Var(C_{t+1} \mid \Lambda_{t+1}, R_{t+1}) &= \rho \Lambda_{t+1} R_{t+1} + \frac{\rho^2 \Lambda_{t+1} R_{t+1}^2}{k}.
\end{align}
应用全方差公式展开 $\Var(C_{t+1} \mid \mathcal{F}_t)$：
\begin{equation}
  \Var(C_{t+1} \mid \mathcal{F}_t) = \E\big[ \Var(C_{t+1} \mid \Lambda_{t+1}, R_{t+1}) \big] + \Var\big[ \E(C_{t+1} \mid \Lambda_{t+1}, R_{t+1}) \big].
  \label{eq:eve-renewal}
\end{equation}
分别计算式~\eqref{eq:eve-renewal} 中的两项：
\begin{enumerate}
  \item \textbf{第一项（条件方差期望）：}根据条件独立性，在历史信息集 $\mathcal{F}_t$ 给定下，未来的环境 $R_{t+1}$ 与历史生成的发病存量 $\Lambda_{t+1}$ 相互独立：
  \begin{align}
    \E\big[ \Var(C_{t+1} \mid \Lambda_{t+1}, R_{t+1}) \big] &= \E\left[ \rho \Lambda_{t+1} R_{t+1} + \frac{\rho^2 \Lambda_{t+1} R_{t+1}^2}{k} \right] \nonumber \\
    &= \rho \E(\Lambda_{t+1}) \E(R_{t+1}) + \frac{\rho^2 \E(\Lambda_{t+1}) \E(R_{t+1}^2)}{k} \nonumber \\
    &= \rho \bar{\Lambda} \Rbar + \frac{\rho^2 \bar{\Lambda}(\Rbar^2 + v_R)}{k} \nonumber \\
    &= m \left[ \Rbar + \frac{\rho(\Rbar^2 + v_R)}{k} \right] = m A.
  \end{align}
  \item \textbf{第二项（条件均值方差）：}计算乘积项的方差：
  \begin{align}
    \Var\big( \rho \Lambda_{t+1} R_{t+1} \big) &= \rho^2 \Var(\Lambda_{t+1} R_{t+1}) \nonumber \\
    &= \rho^2 \Big\{ \E(\Lambda_{t+1}^2 R_{t+1}^2) - [\E(\Lambda_{t+1} R_{t+1})]^2 \Big\} \nonumber \\
    &= \rho^2 \Big\{ \E(\Lambda_{t+1}^2) \E(R_{t+1}^2) - [\E(\Lambda_{t+1})]^2 [\E(R_{t+1})]^2 \Big\} \nonumber \\
    &= \rho^2 \Big\{ [\Var(\Lambda_{t+1}) + \bar{\Lambda}^2](\Rbar^2 + v_R) - \bar{\Lambda}^2 \Rbar^2 \Big\} \nonumber \\
    &= \rho^2 \Big\{ \Var(\Lambda_{t+1})(\Rbar^2 + v_R) + \bar{\Lambda}^2(\Rbar^2 + v_R) - \bar{\Lambda}^2 \Rbar^2 \Big\} \nonumber \\
    &= \rho^2 \Big\{ \bar{\Lambda}^2 v_R + \Var(\Lambda_{t+1})(\Rbar^2 + v_R) \Big\} \nonumber \\
    &= \rho^2 \bar{\Lambda}^2 \left[ v_R + \frac{\Var(\Lambda_{t+1})}{\bar{\Lambda}^2}(\Rbar^2 + v_R) \right] \nonumber \\
    &= m^2 \Big[ v_R + \CV^2(\Lambda_{t+1} \mid \mathcal{F}_t)\,(\Rbar^2 + v_R) \Big].
  \end{align}
\end{enumerate}
将两项相加，两端同除以 $m^2 = [\rho \E(\Lambda_{t+1} \mid \mathcal{F}_t)]^2$：
\begin{equation}
  \frac{\Var(C_{t+1} \mid \mathcal{F}_t)}{m^2} = \frac{m A + m^2 [v_R + \CV^2(\Lambda_{t+1})(\Rbar^2+v_R)]}{m^2} = v_R + \frac{A}{m} + \CV^2(\Lambda_{t+1} \mid \mathcal{F}_t)\,(\Rbar^2 + v_R).
\end{equation}
式~\eqref{eq:info-decomp-der} 获证。
\end{proof}

\begin{insightbox}{物理洞见 12.1：潜变量不确定性带来的结构性抬升}
当真实感染压力已知时（例如前瞻性追踪全部病例），$\CV(\Lambda_{t+1}) = 0$，式~\eqref{eq:info-decomp-der} 精确退化为单代独立前瞻预测界限 $v_R + A/m$。
在被动时间序列中，由于不知道过去到底有多少隐匿感染者在发病，感染压力的相对方差 $\CV^2(\Lambda_{t+1})$ 形成了一个\textbf{纯粹由监测不完全性诱发的外生方差底板}，且随世代间隔的展宽和漏报的加剧而放大。
\end{insightbox}

\subsection{分母二项抽样噪声的二阶 Delta 展开（式 (b-delta) 证明）}

在经验拟合中，由于无法使用潜在期望规模 $m$，回归方程以观测的分母规模 $D_t = \sum_s w_s C_{t+1-s}$ 作为自变量拟合 $L_t = a + b/D_t$。分母包含二项抽样噪声，属于经典计量经济学中的“测量误差自变量（Errors-in-variables）”。以下显式证明分母噪声为何将经验斜率 $b$ 抬高。

\begin{theorembox}{推导 12.2：分母噪声抬高斜率的二阶 Delta 展开证明}
设给定真实感染历史 $\mathcal{H}_t = (I_t, I_{t-1}, \dots)$，潜在理论规模为 $m = \rho \Lambda_{t+1} = \rho \sum_s w_s I_{t+1-s}$。分母相对噪声定义为 $\delta \equiv \frac{D_t - m}{m}$。
在二阶 Delta 展开下，归一化均方损失关于 $1/m$ 的回归斜率近似满足式~\eqref{eq:b-delta}：
\begin{equation}
  b \approx \frac{A + (1 - \rho)(\Rbar^2 + v_R)\sum_{s} w_s^2}{\Rbar^2}.
  \label{eq:b-delta}
\end{equation}
\end{theorembox}

\begin{proof}
根据各周独立二项报告抽样假设，给定感染历史 $\mathcal{H}_t$：
\begin{equation}
  C_{t+1-s} \mid I_{t+1-s} \sim \Bin(I_{t+1-s}, \rho).
\end{equation}
计算分母的条件期望与条件方差：
\begin{align}
  \E(D_t \mid \mathcal{H}_t) &= \E\left( \sum_s w_s C_{t+1-s} \ \Big|\ \mathcal{H}_t \right) = \sum_s w_s (\rho I_{t+1-s}) = \rho \sum_s w_s I_{t+1-s} = m, \\
  \Var(D_t \mid \mathcal{H}_t) &= \Var\left( \sum_s w_s C_{t+1-s} \ \Big|\ \mathcal{H}_t \right) = \sum_s w_s^2 \Var(C_{t+1-s} \mid I_{t+1-s}) \nonumber \\
  &= \sum_s w_s^2 \big[ \rho (1 - \rho) I_{t+1-s} \big] = \rho (1 - \rho) \sum_s w_s^2 I_{t+1-s}.
\end{align}
在局部平衡阶段，$I_{t+1-s} \approx \Lambda_{t+1} = m/\rho$。代入得：
\begin{equation}
  \Var(D_t \mid \mathcal{H}_t) \approx \rho (1 - \rho) \frac{m}{\rho} \sum_s w_s^2 = m (1 - \rho) \sum_s w_s^2 \equiv m c_w,
\end{equation}
其中记常数 $c_w \equiv (1 - \rho)\sum_s w_s^2$。由此计算相对扰动 $\delta = \frac{D_t - m}{m}$ 的统计矩：
\begin{equation}
  \E(\delta \mid \mathcal{H}_t) = 0, \qquad \E(\delta^2 \mid \mathcal{H}_t) = \frac{\Var(D_t \mid \mathcal{H}_t)}{m^2} \approx \frac{m c_w}{m^2} = \frac{c_w}{m}.
\end{equation}
增长统计量为 $G_t = \frac{C_{t+1}}{D_t} = \frac{C_{t+1}}{m(1+\delta)} = X (1 + \delta)^{-1}$，其中定义无噪分母统计量 $X \equiv C_{t+1}/m$。
由未来抽样与历史抽样的独立性，在给定 $\mathcal{H}_t$ 下，$X$ 与 $\delta$ \textbf{相互独立}。
已知无噪统计量的条件矩为：$\E(X) = \Rbar$，$\Var(X) = v_R + A/m$，从而 $\E(X^2) = \Var(X) + [\E(X)]^2 = v_R + \Rbar^2 + A/m$。

将逆函数展开到 $\delta$ 的二次项：
\begin{equation}
  (1 + \delta)^{-1} = 1 - \delta + \delta^2 - \delta^3 + O(\delta^4).
\end{equation}
利用独立性分别求条件期望与二阶矩：
\begin{align}
  \E(G_t \mid \mathcal{H}_t) &= \E\big[ X (1 - \delta + \delta^2 + \dots) \big] = \E(X) \big[ 1 - \E(\delta) + \E(\delta^2) + \dots \big] \nonumber \\
  &= \Rbar \left( 1 + \frac{c_w}{m} \right) + O(m^{-2}).
\end{align}
计算平方项期望 $\E(G_t^2 \mid \mathcal{H}_t)$：
\begin{align}
  \E(G_t^2 \mid \mathcal{H}_t) &= \E\big[ X^2 (1 + \delta)^{-2} \big] = \E(X^2) \, \E\big[ 1 - 2\delta + 3\delta^2 + O(\delta^3) \big] \nonumber \\
  &= \E(X^2) \left( 1 + 3\frac{c_w}{m} \right) + O(m^{-2}) \nonumber \\
  &= \left( \Rbar^2 + v_R + \frac{A}{m} \right) \left( 1 + 3\frac{c_w}{m} \right) + O(m^{-2}) \nonumber \\
  &= \Rbar^2 + v_R + \frac{A}{m} + 3(\Rbar^2 + v_R)\frac{c_w}{m} + O(m^{-2}).
\end{align}
进而计算方差：
\begin{align}
  \Var(G_t \mid \mathcal{H}_t) &= \E(G_t^2) - [\E(G_t)]^2 \nonumber \\
  &= \left[ \Rbar^2 + v_R + \frac{A}{m} + 3(\Rbar^2 + v_R)\frac{c_w}{m} \right] - \Rbar^2 \left( 1 + \frac{c_w}{m} \right)^2 + O(m^{-2}) \nonumber \\
  &= \Rbar^2 + v_R + \frac{A}{m} + 3(\Rbar^2 + v_R)\frac{c_w}{m} - \Rbar^2 \left( 1 + 2\frac{c_w}{m} \right) + O(m^{-2}) \nonumber \\
  &= v_R + \frac{A}{m} + 3(\Rbar^2 + v_R)\frac{c_w}{m} - 2\Rbar^2\frac{c_w}{m} + O(m^{-2}) \nonumber \\
  &= v_R + \frac{A + (\Rbar^2 + 3v_R)c_w}{m} + O(m^{-2}).
  \label{eq:var-Gt-raw}
\end{align}
注意式~\eqref{eq:var-Gt-raw} 中包含项 $(\Rbar^2 + 3v_R)c_w$。
现在将方差用条件均值的平方 $[\E(G_t \mid \mathcal{H}_t)]^2 = \Rbar^2(1 + 2c_w/m) + O(m^{-2})$ 进行标准化归一：
\begin{align}
  \frac{\Var(G_t \mid \mathcal{H}_t)}{[\E(G_t \mid \mathcal{H}_t)]^2} &= \frac{v_R + \frac{A + (\Rbar^2 + 3v_R)c_w}{m}}{\Rbar^2 \left( 1 + \frac{2c_w}{m} \right)} + O(m^{-2}) \nonumber \\
  &= \frac{1}{\Rbar^2} \left[ v_R + \frac{A + (\Rbar^2 + 3v_R)c_w}{m} \right] \left( 1 - \frac{2c_w}{m} \right) + O(m^{-2}) \nonumber \\
  &= \frac{1}{\Rbar^2} \left[ v_R - v_R \frac{2c_w}{m} + \frac{A + (\Rbar^2 + 3v_R)c_w}{m} \right] + O(m^{-2}) \nonumber \\
  &= \frac{v_R}{\Rbar^2} + \frac{A + (\Rbar^2 + 3v_R)c_w - 2v_R c_w}{\Rbar^2 m} + O(m^{-2}) \nonumber \\
  &= \frac{v_R}{\Rbar^2} + \frac{A + (\Rbar^2 + v_R)c_w}{\Rbar^2 m} + O(m^{-2}).
\end{align}
注意到极为精巧的一幕：$-2v_R c_w$ 恰好将 $3v_R$ 消减为 $v_R$！
将 $c_w = (1 - \rho)\sum_s w_s^2$ 代回分子中的 $1/m$ 系数，即得出理论斜率修正项：
\begin{equation}
  b \approx \frac{A + (1 - \rho)(\Rbar^2 + v_R)\sum_s w_s^2}{\Rbar^2}.
\end{equation}
推导完毕。
\end{proof}

\section{条件有偏估计：线性收缩与经验贝叶斯}
\label{sec:part13}

\subsection{突破条件无偏下确界的理论可能性}

在定理 4.1 中，下确界 $v_R$ 的刚性约束建立在\textbf{条件无偏性}（$\E(\tilde{\theta} \mid R_t) = R_t$）这一前提假设之上。在现代统计学与机器推断中，诸如 James--Stein 收缩估计量、经验贝叶斯（Empirical Bayes）以及状态空间平滑算法（如 EpiFilter，\citet{parag2021improved}）允许引入确定性条件偏差，以期换取方差的大幅缩减。

\subsection{线性收缩估计量的均方误差闭式推导}

\begin{derivationbox}{推导 13.1：最优收缩强度与先验偏差容限}
构造经典凸组合线性收缩估计量：
\begin{equation}
  \hat{\theta}_\lambda \equiv (1 - \lambda) \frac{C_t}{m} + \lambda \mu_0, \qquad \lambda \in [0, 1],
  \label{eq:shrink-def}
\end{equation}
其中 $\mu_0$ 为先验锚定中心（如多季历史基准或流行病学常模），$\lambda$ 为收缩强度。记先验中心与真实环境均值的偏差为 $\delta \equiv \mu_0 - \Rbar$，无偏估计量方差为 $V \equiv v_R + A/m$。
计算总均方误差：
\begin{align}
  \hat{\theta}_\lambda - \Rbar &= (1 - \lambda)\left( \frac{C_t}{m} - \Rbar \right) + \lambda(\mu_0 - \Rbar) \nonumber \\
  &= (1 - \lambda)\left( \frac{C_t}{m} - \Rbar \right) + \lambda \delta.
\end{align}
由于 $\E(C_t/m - \Rbar) = 0$，平方求期望时交叉项消失：
\begin{equation}
  \text{MSE}(\lambda) \equiv \E\big[(\hat{\theta}_\lambda - \Rbar)^2\big] = (1 - \lambda)^2 V + \lambda^2 \delta^2.
\end{equation}
对收缩参数 $\lambda$ 求一阶导数寻找极值：
\begin{equation}
  \frac{\mathrm{d}\,\text{MSE}(\lambda)}{\mathrm{d}\lambda} = -2(1 - \lambda)V + 2\lambda \delta^2 = 0 \implies \lambda(V + \delta^2) = V.
\end{equation}
由此得出最优收缩权重（Oracle shrinkage intensity）：
\begin{equation}
  \boxed{\, \lambda^* = \frac{V}{V + \delta^2} \,}
  \label{eq:opt-lambda}
\end{equation}
将最优权重代回均方误差公式：
\begin{align}
  \text{MSE}(\lambda^*) &= \left( 1 - \frac{V}{V+\delta^2} \right)^2 V + \left( \frac{V}{V+\delta^2} \right)^2 \delta^2 \nonumber \\
  &= \left( \frac{\delta^2}{V+\delta^2} \right)^2 V + \frac{V^2 \delta^2}{(V+\delta^2)^2} \nonumber \\
  &= \frac{V \delta^4 + V^2 \delta^2}{(V+\delta^2)^2} = \frac{V \delta^2 (V + \delta^2)}{(V+\delta^2)^2} = \frac{V \delta^2}{V + \delta^2}.
  \label{eq:opt-mse-shrink}
\end{align}
\end{derivationbox}

\begin{theorembox}{定理 13.1（收缩估计击穿环境方差的充要条件）}
最优收缩估计量的均方误差严格小于不可约环境方差 $v_R$（即 $\text{MSE}(\lambda^*) < v_R$）当且仅当先验偏差满足：
\begin{equation}
  \delta^2 < v_R \left( 1 + \frac{m}{\mx} \right).
  \label{eq:shrink-barrier}
\end{equation}
\end{theorembox}

\begin{proof}
根据式~\eqref{eq:opt-mse-shrink}，求解不等式：
\begin{equation}
  \frac{V \delta^2}{V + \delta^2} < v_R \iff V \delta^2 < v_R V + v_R \delta^2 \iff (V - v_R) \delta^2 < v_R V.
\end{equation}
已知 $V - v_R = A/m$，且 $V = v_R + A/m$。代入得：
\begin{equation}
  \frac{A}{m} \delta^2 < v_R \left( v_R + \frac{A}{m} \right).
\end{equation}
两侧同乘以 $\frac{m}{A}$：
\begin{equation}
  \delta^2 < v_R \left( \frac{v_R m}{A} + 1 \right).
\end{equation}
利用方差平衡规模定义 $\mx = A / v_R \implies \frac{v_R}{A} = \frac{1}{\mx}$，立即得出：
\begin{equation}
  \delta^2 < v_R \left( 1 + \frac{m}{\mx} \right).
\end{equation}
定理获证。
\end{proof}

\begin{insightbox}{权衡与警示：收缩方差红利的代价}
式~\eqref{eq:shrink-barrier} 证明：只要先验中心 $\mu_0$ 与总体均值的偏差平方 $\delta^2$ 落在容许椭球内，经验贝叶斯确实可以战胜环境方差 $v_R$。
在基准参数（$m = 50, v_R = 0.04$）下，只要先验偏差 $|\delta| < 0.279$，收缩估计的 MSE 就低于 $0.04$；当偏差仅为 $0.1$ 时，MSE 甚至只有 $0.0089$（仅为 $v_R$ 的 22\%）。

\textbf{但天上没有免费的午餐。其惨痛代价是引入了状态依赖的条件偏差：}
\begin{equation}
  \text{Bias}(\hat{\theta}_\lambda \mid R_t) \equiv \E(\hat{\theta}_\lambda \mid R_t) - R_t = (1 - \lambda) R_t + \lambda \mu_0 - R_t = \lambda(\mu_0 - R_t).
\end{equation}
条件偏差随当期传播环境 $R_t$ 偏离常模的幅度\textbf{线性暴增}！在常态稳定期（$R_t \approx \Rbar$），估计量非常精准；而在新发突变株引发超级暴发或断崖式骤降的非常态危机时刻（$|R_t - \mu_0| \gg 0$），收缩估计量因强烈被锚定在历史中心，发生灾难性的钝化与误判。这恰恰违背了公共卫生危机干预的初衷。
\end{insightbox}

\section{与 Parag 等信息论框架的深度比较}
\label{sec:part14}

\subsection{Parag 等经典泊松更新信息论回顾}

近年来，牛津大学的 Kris Parag 与 Christl Donnelly \citep{parag2022quantifying, parag2022fundamental} 将经典信息论与统计物理引入传染病再生数估计，建立了以 Fisher 信息量与互信息（Mutual Information）为核心的评估框架。其基准模型设定如下：
\begin{enumerate}
  \item 发病计数服从泊松分布：$I_t \mid \Lambda_t, R_t \sim \text{Poisson}(R_t \Lambda_t)$；
  \item 其条件对数似然为 $\ell(R_t) = I_t \ln(R_t \Lambda_t) - R_t \Lambda_t - \ln(I_t!)$；
  \item 经典泊松 Fisher 信息量为：
  \begin{equation}
    \mathcal{I}_{\text{Poisson}}(R_t) \equiv -\E\left( \frac{\partial^2 \ell}{\partial R_t^2} \right) = \frac{\E(I_t)}{R_t^2} = \frac{R_t \Lambda_t}{R_t^2} = \frac{\Lambda_t}{R_t}.
  \end{equation}
  \item 相对 Fisher 信息量（即无量纲标准化信息）：
  \begin{equation}
    \mathcal{F}(R_t) \equiv R_t^2 \, \mathcal{I}_{\text{Poisson}}(R_t) = R_t \Lambda_t = \E(I_t).
  \end{equation}
  这给出了极具穿透力的经典结论：\textbf{再生数的信息量在数值上恰好等于当期累积的期望病例数}。
\end{enumerate}

\subsection{四项核心数学连接与区别定理}

\begin{theorembox}{推导 14.1：向负二项超传播与环境随机性拓展的四大理论定理}
将本文的微观负二项分支与双目标框架与 Parag 的泊松框架对齐，可严格证明以下四项数学结论：
\end{theorembox}

\begin{enumerate}[label=\textbf{结论 \arabic*：}]
  \item \textbf{渐近退化性：无超传播且无环境波动时严格重合。}
  在本文模型中，令超传播参数 $k \to \infty$（无超级传播）、报告率 $\rho = 1$（全量报告）、祖代规模记为 $\Lambda_t$。根据式~\eqref{eq:fisher-def}：
  \begin{equation}
    \lim_{k \to \infty, \rho \to 1} \mathcal{I}(R_t) = \lim_{k \to \infty, \rho \to 1} \frac{n}{\frac{R_t}{\rho} + \frac{R_t^2}{k}} = \frac{n}{R_t + 0} = \frac{\Lambda_t}{R_t} = \mathcal{I}_{\text{Poisson}}(R_t).
  \end{equation}
  并且由于 $v_R = 0$，目标 1 与目标 2 重合，$A \to \Rbar = R_t$，$V(m) = 0 + \frac{R_t}{n} = \frac{1}{\mathcal{I}_{\text{Poisson}}}$。本文结论完整包含了 Parag 的经典结论作为特例。

  \item \textbf{超传播缓冲效应：不完全报告在负二项下的信息衰减小于泊松情形。}
  在泊松模型下，若报告率从 $1$ 降为 $\rho$，观测方差从 $R_t/n$ 升为 $R_t/(\rho n)$，相对信息损失比例严格等于 $1 - \rho$：
  \begin{equation}
    \frac{\mathcal{I}_{\rho=1} - \mathcal{I}_\rho}{\mathcal{I}_{\rho=1}} = \frac{\frac{n}{R_t} - \frac{\rho n}{R_t}}{\frac{n}{R_t}} = 1 - \rho.
  \end{equation}
  但在超传播（$k < \infty$）的真实世界下，代入式~\eqref{eq:fisher-def}：
  \begin{align}
    \frac{\mathcal{I}_{\rho=1} - \mathcal{I}_\rho}{\mathcal{I}_{\rho=1}} &= 1 - \frac{\frac{n}{R_t/\rho + R_t^2/k}}{\frac{n}{R_t + R_t^2/k}} = 1 - \frac{R_t + \frac{R_t^2}{k}}{\frac{R_t}{\rho} + \frac{R_t^2}{k}} \nonumber \\
    &= \frac{\left(\frac{R_t}{\rho} + \frac{R_t^2}{k}\right) - \left(R_t + \frac{R_t^2}{k}\right)}{\frac{R_t}{\rho} + \frac{R_t^2}{k}} = \frac{R_t\left(\frac{1}{\rho} - 1\right)}{\frac{R_t}{\rho}\left(1 + \frac{\rho R_t}{k}\right)} \nonumber \\
    &= \frac{1 - \rho}{1 + \frac{\rho R_t}{k}} < 1 - \rho.
  \end{align}
  由于分母中存在严格正项 $\frac{\rho R_t}{k} > 0$，超传播使得信息损失比例严格低于 $1 - \rho$！
  \textbf{直观解释：}超级传播个体造成的传染力方差本身已占主导，后置的二项抽样漏报只是在已经严重离散的数据上追加了一层次要随机性，因而边际信息破坏力相对减弱。

  \item \textbf{方差稳定变换在负二项下不可行。}
  在泊松模型中，Parag 利用 Anscombe 平方根变换 $Z_t = 2\sqrt{I_t}$，其渐近方差由一阶泰勒展开给出：
  \begin{equation}
    \Var(2\sqrt{I_t}) \approx \left[ 2 \cdot \frac{1}{2\sqrt{\E I_t}} \right]^2 \Var(I_t) = \frac{1}{R_t \Lambda_t} \cdot (R_t \Lambda_t) \equiv 1.
  \end{equation}
  由此彻底消除了方差对未知参数 $R_t$ 的依赖性。
  但在负二项超传播下，方差为 $\E(C_t) + \frac{\rho^2 n R_t^2}{k} = m R_t + \frac{\rho m R_t^2}{k}$。由于含有均值的二次项，任何幂函数变换均无法消去 $R_t$。

  \item \textbf{统计信息与物理离散度的根本边界。}
  Parag 的信息论体系是“\textbf{推断视角（Inferential perspective）}”：它回答“给定确定的现实，怎样从有限观测中最大限度抽取信息”。因此只要样本无限大，后验分布必然无限坍缩为狄拉克函数（Dirac delta），方差必然为零。
  而本文的环境随机性是“\textbf{生成视角（Generative perspective）}”：环境方差 $v_R$ 是大自然在每一世代施加的物理扰动。即便观测没有任何噪声且样本无限大，明天的传播强度相对于过去的均值依然具有物理发散度 $v_R$。这是生态学生物学现实，无法仅通过改进算法或堆积病例来消除。
\end{enumerate}

\section{理论答辩与核心追问指南}
\label{sec:part15}

为了在学术报告、学位论文答辩或同行评议中能够清晰、准确、有理有据地回应质疑，本节整理了五大最关键的追问与标准应答方案。

\subsection{追问 1：大样本下根据大数定律方差必趋于 0，你的下界 $v_R$ 是否违背了中心极限定理？}

\begin{derivationbox}{标准答辩要点}
\textbf{核心回应：完全不违背。关键在于严格区分“估计目标 1”与“估计目标 2”。}
\begin{enumerate}
  \item 如果估计目标是\textbf{当期的条件瞬时值 $R_t$}（估计目标 1），在给定当期环境时，根据中心极限定理与大数定律，自然估计量 $\widehat{R}_t = C_t/m$ 的方差为：
  \begin{equation}
    \Var(\widehat{R}_t \mid R_t) = \frac{R_t}{m} + \frac{\rho R_t^2}{km} \xrightarrow{m \to \infty} 0.
  \end{equation}
  它确实严格单调收敛于 0，并达到了 Cramér--Rao 下界，完全符合经典大样本理论。
  \item 但如果估计目标是\textbf{跨期平均再生数 $\Rbar$}（估计目标 2），估计误差由两部分复合而成：
  \begin{equation}
    \widehat{R}_t - \Rbar = \underbrace{(\widehat{R}_t - R_t)}_{\to 0} + \underbrace{(R_t - \Rbar)}_{\text{不随 } m \text{ 改变}}.
  \end{equation}
  第一部分是微观抽样误差，随规模 $m \to \infty$ 确切收敛至 0；但第二部分是客观自然界在每一世代发生的物理波动。因为我们只是在\textbf{单期内}增加了病例样本量，并没有跨越时间采集更多独立的天气与人群接触样本，所以 $R_t$ 本身相对于长期均值的偏离不会因为你在这一周内测了 1000 人还是 100 万人而消失。全方差公式给出的正是 $v_R + A/m \to v_R$。
\end{enumerate}
\end{derivationbox}

\subsection{追问 2：你的结论是否意味着“增加监测病例规模没有用处”？}

\begin{derivationbox}{标准答辩要点}
\textbf{核心回应：绝非没有用处。方差函数 $V(m)$ 处处严格单调递减，扩大规模始终能压低方差，但其边际效率随规模急剧衰减。}
\begin{enumerate}
  \item 方差函数的一阶导数处处为负：$V'(m) = -A/m^2 < 0$。因此增加病例数始终是有益的。
  \item 关键在于\textbf{性价比与边际替代}：
  \begin{itemize}
    \item 当规模较小（$m \ll \mx$）时，总方差中抽样项占 90\% 以上，此时扩大病例规模是最核心的抓手，翻倍扩样可压制近 50\% 的方差；
    \item 当规模跨越平衡点（$m > \mx$）后，可约抽样项在总方差中已降至 50\% 以下（当规模进一步达到 $m > 4\mx$ 时降至 20\% 以下）。此时若继续将经费用于单期扩样，面临显著钝化的边际回报；
    \item 此时应当将资源\textbf{边际转向增加独立观测周期数 $T$、跨区域空间汇总 $J$ 或统一报告口径 $\tau \to 0$}，因为这些杠杆直接作用于单期扩样无法消除的 $v_R$ 与 $v_H$。
  \end{itemize}
\end{enumerate}
\end{derivationbox}

\subsection{追问 3：为什么美国流感实证截距 $a = 0.0991$ 显著大于理论单代情景 $v_R/\Rbar^2 \approx 0.014$--$0.028$？}

\begin{derivationbox}{标准答辩要点}
\textbf{核心回应：实证拟合截距捕捉的是经验模型中未被规模项解释的全部非抽样误差，由三方面机制共同放大：}
\begin{enumerate}
  \item \textbf{世代聚合与潜变量不确定性：}流感世代间隔约为 3 天，周度监测数据聚合了约两代感染。根据第~\ref{sec:part12} 节定理 12.1，时间序列更新过程引入了潜在感染压力的相对方差 $\CV^2(\Lambda_{t+1}) (\Rbar^2 + v_R)$，将理论截距抬升数倍；
  \item \textbf{未观测的微观报告波动：}NHSN 医院填报比例代理仅能记录“医院是否提交数据”，无法记录已填报医院内部由于就诊门槛、检测盒短缺造成的微观漏报率波动 $\tau$。根据式~\eqref{eq:vH-def}，任何未被代理变量捕捉的残余报告异质性 $v_H = \tau^2(\Rbar^2+v_R)$ 都会机械并入经验截距；
  \item \textbf{审慎定性：}正文因此严格遵循审慎原则，明确指出“拟合截距反映该经验模型中未被规模项解释的误差成分，不能将其唯一、简单地归因于纯生物学环境方差”。
\end{enumerate}
\end{derivationbox}

\subsection{追问 4：平衡规模 $\mx = 53.1$ 例是否意味着“一个辖区每周只要采 53 例就足够了”？}

\begin{derivationbox}{标准答辩要点}
\textbf{核心回应：完全不是。平衡规模是比较基准，而非满足特定性能的操作阈值。}
\begin{enumerate}
  \item $\mx$ 的唯一定义是抽样项等于环境项的方差切分点（$\frac{A}{\mx} = v_R$）。
  \item 监测系统需要多少病例，必须\textbf{严格取决于任务类型}：
  \begin{itemize}
    \item 如果任务是\textbf{长期趋势与基准评估（目标 2）}：当达到 53 例左右时，抽样误差已经与环境背景噪音同量级，继续单点扩样的边际收益已经大幅收窄；
    \item 如果任务是\textbf{暴发早预警（目标 1）}：在小样本离散情境下，若要以 80\% 的统计功效可靠检出 15\% 的疫情回升，根据第~\ref{sec:part11} 节推导，必须稳定捕获 $m \approx 527$ 例（是 $\mx$ 的 10 倍）。
  \end{itemize}
  因此，监测设计必须将预警哨点（追求大样本、高敏感度）与趋势评估哨点进行结构性解耦。
\end{enumerate}
\end{derivationbox}

\subsection{追问 5：在时间正相关环境（$\phi > 0$）下，为什么预算小时要做大单期，预算大时要做多期？}

\begin{derivationbox}{标准答辩要点}
\textbf{核心回应：这是信息边际冗余与抽样方差非线性爆炸之间的精妙权衡。}
\begin{enumerate}
  \item \textbf{当总经费很低（$\mathcal{B} < \mathcal{B}^\dagger$）时：}若教条地将微薄的资金切分成许多期，会导致每一期的样本量极其微小（如每周只有 2--3 例）。由于抽样方差按 $1/m$ 呈双曲线爆炸，极度匮乏的单期样本会造成巨大的离散噪声，甚至频繁出现零病例导致估计崩溃。此时将资金聚合在少数几期、把单期做实做大，边际回报最高；
  \item \textbf{当总经费充裕（$\mathcal{B} > \mathcal{B}^\dagger$）时：}单期样本早已越过平衡规模，继续扩大单期规模所能削减的方差已极其有限。而环境存在正自相关意味着相邻周提供的信息高度重叠。此时最优策略必须是拉长时间跨度、增加采样期数 $T$，尽可能收集跨越不同自相关周期的独立环境扰动样本。
\end{enumerate}
\end{derivationbox}


\newpage
\bibliographystyle{plainnat}
\small
\setlength{\bibsep}{2pt}
\bibliography{bibliography}

\end{document}
